% Generated mechanically from frozen input canon/ko-kr/integration-20260908-r10-p03-register/inputs/p03/ko/dpa.tex.
% Do not edit this generated file; edit the bound locale source instead.

% OK, start here.
%

\setcounter{chapter}{22}
\chapter{분할 거듭제곱 대수}



\phantomsection
\label{dpa-section-phantom}


\section{서론}
\label{dpa-section-introduction}

\noindent
이 장에서는 분할 거듭제곱 대수와 그것으로 할 수 있는 일을
논한다. 참고문헌은 \cite{Berthelot}이다.





\section{분할 거듭제곱}
\label{dpa-section-divided-powers}

\noindent
이 절에서는 분할 거듭제곱 환에 관한 몇 가지 결과를 모은다.
빈 곱은 $1$이어야 하므로 $0! = 1$이라는 관례를 사용한다.

\begin{definition}
\label{dpa-definition-divided-powers}
$A$를 환이라 하고 $I$를 $A$의 아이디얼이라 하자. 사상들의 모음
$\gamma_n : I \to I$, $n > 0$을 $I$ 위의
{\it 분할 거듭제곱 구조}라 함은 모든 $n \geq 0$, $m > 0$,
$x, y \in I$, $a \in A$에 대해 다음이 성립한다는 뜻이다.
\begin{enumerate}
\item $\gamma_1(x) = x$이고, 또한 $\gamma_0(x) = 1$로 놓는다.
\item $\gamma_n(x)\gamma_m(x) = \frac{(n + m)!}{n! m!} \gamma_{n + m}(x)$이다.
\item $\gamma_n(ax) = a^n \gamma_n(x)$이다.
\item $\gamma_n(x + y) = \sum_{i = 0, \ldots, n} \gamma_i(x)\gamma_{n - i}(y)$이다.
\item $\gamma_n(\gamma_m(x)) = \frac{(nm)!}{n! (m!)^n} \gamma_{nm}(x)$이다.
\end{enumerate}
\end{definition}

\noindent
정의에 나타나는 유리수 $\frac{(n + m)!}{n! m!}$과
$\frac{(nm)!}{n! (m!)^n}$은 사실 정수이다. 첫째 것은
$n + m$개 가운데 $n$개를 고르는 방법의 수이고, 둘째 것은
$nm$개의 대상들로 이루어진 한 무리를 $m$개씩인 $n$개의
무리로 나누는 방법의 수이다. 다음은 $\gamma_n(x)$가 $I$에서
$x^n/n!$을 대신한다는 것을 보여 주는 정의에 관한 몇 가지
관찰이다.

\begin{lemma}
\label{dpa-lemma-silly}
$A$를 환이라 하고 $I$를 $A$의 아이디얼이라 하자.
\begin{enumerate}
\item $\gamma$가 $I$ 위의 분할 거듭제곱 구조이면\footnote{여기와
이후에서 $\gamma$는 $I$에서 $I$로 가는 사상들의 열
$\gamma_1, \gamma_2, \gamma_3, \ldots$을 줄여 쓴 것이다.},
$n \geq 1$, $x \in I$에 대해
$n! \gamma_n(x) = x^n$이다.
\end{enumerate}
$A$가 $\mathbf{Z}$-가군으로서 비틀림이 없다고 가정하자.
\begin{enumerate}
\item[(2)] $I$ 위의 분할 거듭제곱 구조가 존재하면 유일하다.
\item[(3)] $\gamma_n : I \to I$들이 사상들이면
$$
\gamma\text{ is a divided power structure}
\Leftrightarrow
n! \gamma_n(x) = x^n\ \forall x \in I, n \geq 1.
$$
\item[(4)] 아이디얼 $I$가 분할 거듭제곱 구조를 갖기 위한
필요충분조건은, 모든 $n \geq 1$에 대해 $x_i^n \in (n!)I$인
아이디얼로서의 $I$의 생성자 집합 $x_i$가 존재하는 것이다.
\end{enumerate}
\end{lemma}

\begin{proof}
(1)의 증명. $\gamma$가 분할 거듭제곱 구조이면, 조건 (2)를
$n$과 $m$ 대신 $1$과 $n-1$에 적용하여
$n \gamma_n(x) = \gamma_1(x)\gamma_{n - 1}(x)$를 얻는다.
따라서 귀납법과 조건 (1)에 의해 $n! \gamma_n(x) = x^n$이다.

\medskip\noindent
$A$가 $\mathbf{Z}$-가군으로서 비틀림이 없다고 가정하자.
(2)의 증명. 이는 (1)에서 자명하다.

\medskip\noindent
(3)의 증명. 모든 $x \in I$와 $n \geq 1$에 대해
$n! \gamma_n(x) = x^n$이라고 가정하자.
$A \subset A \otimes_{\mathbf{Z}} \mathbf{Q}$이므로
$A$가 $\mathbf{Q}$-대수인 경우에 Definition
\ref{dpa-definition-divided-powers}의 공리 (1)--(5)를 증명하면
충분하다. 이 경우 $\gamma_n(x) = x^n/n!$이고 (1)--(5)는
곧바로 확인된다. 예를 들어 (4)는 이항공식
$$
(x + y)^n = \sum_{i = 0, \ldots, n} \frac{n!}{i!(n - i)!} x^iy^{n - i}
$$
에 대응한다. 계수들이 올바른지 확인하기 위해 독자가 직접
이 검산을 해 보기를 권한다.

\medskip\noindent
(4)의 증명. 모든 $n \geq 1$에 대해 $x_i^n \in (n!)I$인
아이디얼로서의 $I$의 생성자 $x_i$가 주어졌다고 하자.
모든 $x \in I$에 대해 $x^n \in (n!)I$라고 주장한다.
이 주장이 성립하면 $\gamma_n(x) = x^n/n!$로 놓을 수 있고,
이는 (3)에 의해 분할 거듭제곱 구조이다.
주장을 증명하기 위해 먼저 $x = ax_i$일 때 성립함을 관찰한다.
따라서 이 주장은 아벨군으로서의 $I$의 생성자 집합에서 성립한다.
그 생성자들로 쓴 식의 길이에 관한 귀납법에 의해, $x$와 $y$에
대해 주장이 성립할 때 $x + y$에 대해서도 성립함을 보이면
충분하다. 이는 이항정리에서 즉시 따른다.
\end{proof}

\begin{example}
\label{dpa-example-ideal-generated-by-p}
$p$를 소수라 하자. $p$로 나누어지지 않는 모든 정수 $n$이
가역인 환 $A$, 즉 $\mathbf{Z}_{(p)}$-대수 $A$를 생각하자.
그러면 $I = pA$는 표준적인 분할 거듭제곱 구조를 갖는다.
실제로 $x = pa \in I$가 주어지면
$$
\gamma_n(x) = \frac{p^n}{n!} a^n
$$
으로 놓는다. $n \geq 1$에 대해
$p^n/n! \in p\mathbf{Z}_{(p)}$임을 독자는 즉시 확인할 수 있다
(예를 들어 $n!$의 소인수분해에서 $p$의 지수가
$\left\lfloor n/p \right\rfloor + \left\lfloor n/p^2 \right\rfloor
+ \left\lfloor n/p^3 \right\rfloor + \ldots$라는 사실에서 이를
얻을 수 있다). 따라서 이 정의는 의미가 있으며 사상들의 열
$\gamma_n : I \to I$를 준다. Definition
\ref{dpa-definition-divided-powers}의 조건 (1)--(5)가 성립한다는
것은 곧바로 확인할 수 있는 연습문제이다.
다른 방법으로는 $A_0 = \mathbf{Z}_{(p)}$에 대해 이 정의가
성립함이 자명하고, 그 결과가 Lemma
\ref{dpa-lemma-gamma-extends}에서 따름을 사용할 수 있다.
\end{example}

\noindent
$A$의 임의의 아이디얼 $I$와 $I$ 위의 임의의 분할 거듭제곱
구조 $\gamma$에 대해 $\gamma_n\left(0\right) = 0$임을 관찰한다.
(이는 Definition \ref{dpa-definition-divided-powers}의 공리 (3)에
$a=0$을 적용하면 따른다.)

\begin{lemma}
\label{dpa-lemma-check-on-generators}
$A$를 환, $I$를 $A$의 아이디얼이라 하자.
$\gamma_n : I \to I$, $n \geq 1$을 사상들의 열이라 하자.
다음을 가정하자.
\begin{enumerate}
\item[(a)] 모든 $x, y \in I$에 대해 Definition
\ref{dpa-definition-divided-powers}의 (1), (3), (4)가 성립한다.
\item[(b)] 아이디얼로서의 $I$의 어떤 생성자 집합에 속하는
$x$에 대해 성질 (2)와 (5)가 성립한다.
\end{enumerate}
그러면 $\gamma$는 $I$ 위의 분할 거듭제곱 구조이다.
\end{lemma}

\begin{proof}
이 증명에서 번호 (1), (2), (3), (4), (5)는 Definition
\ref{dpa-definition-divided-powers}에 열거된 조건들을 가리킨다.
(3)을 적용하면, (2)와 (5)가 $x$에 대해 성립할 때 모든
$a \in A$에 대해 $ax$에 대해서도 성립한다.
따라서 (b)는 (2)와 (5)가 아벨군으로서의 $I$의 생성자 집합에서
성립함을 함의한다. 이제 $x, y \in I$라 하자. 그 생성자들로 쓴
식의 길이에 관한 귀납법에 의해, (2)와 (5)가 $x$와 $y$에 대해 성립한다고 할 때
$x + y$에 대해서도 성립함을 보이면 충분하다.

\medskip\noindent
$x + y$에 대한 (2)의 증명. (4)에 의해
$$
\gamma_n(x + y)\gamma_m(x + y) =
\sum\nolimits_{i + j = n,\ k + l = m}
\gamma_i(x)\gamma_k(x)\gamma_j(y)\gamma_l(y)
$$
이다. $x$와 $y$에 대해 (2)를 사용하면 이는
$$
\sum \frac{(i + k)!}{i!k!}\frac{(j + l)!}{j!l!}
\gamma_{i + k}(x)\gamma_{j + l}(y)
$$
와 같다. 이를 전개식
$$
\gamma_{n + m}(x + y) = \sum \gamma_a(x)\gamma_b(y)
$$
와 비교하면, $a + b = n + m$이 주어졌을 때
$$
\sum\nolimits_{i + k = a,\ j + l = b,\ i + j = n,\ k + l = m}
\frac{(i + k)!}{i!k!}\frac{(j + l)!}{j!l!}
=
\frac{(n + m)!}{n!m!}.
$$
임을 증명해야 한다. 이를 직접 논증하는 대신, 다항식환
$\mathbf{Q}[x, y]$의 아이디얼 $I = (x, y)$에 대해서는
$\gamma_n(f) = f^n/n!$, $f \in I$가 $I$ 위의 분할
거듭제곱 구조를 정의하므로 결과가 성립함을 관찰한다.
따라서 위 유리수들의 등식은 참이다.

\medskip\noindent
(1)--(4)가 성립하고 (5)가 $x$와 $y$에 대해 성립한다고 할 때
$x + y$에 대한 (5)의 증명. 다시 유리수들의 등식으로 환원한다.
실제로 (4)를 사용하여
$\gamma_n(\gamma_m(x + y)) = \gamma_n(\sum \gamma_i(x)\gamma_j(y))$로
쓸 수 있다. (4)를 사용하면 $\gamma_n(\gamma_m(x + y))$를
$\gamma_k(\gamma_i(x)\gamma_j(y))$ 꼴 인자들의 곱인 항들의
합으로 쓸 수 있다. $i > 0$이면
\begin{align*}
\gamma_k(\gamma_i(x)\gamma_j(y)) & =
\gamma_j(y)^k\gamma_k(\gamma_i(x)) \\
& = \frac{(ki)!}{k!(i!)^k} \gamma_j(y)^k \gamma_{ki}(x) \\
& =
\frac{(ki)!}{k!(i!)^k} \frac{(kj)!}{(j!)^k} \gamma_{ki}(x) \gamma_{kj}(y)
\end{align*}
이다. 첫째 등식에서는 (3), 둘째 등식에서는 $x$에 대한 (5),
셋째 등식에서는 (2)를 정확히 $k$번 사용했다.
$y$에 대한 (5)를 사용하면 $i = 0$일 때도 같은 등식이 성립한다.
이와 같이 (5)를 제외한 모든 공리를 계속 사용하면
$$
\gamma_n(\gamma_m(x + y)) =
\sum\nolimits_{i + j = nm} c_{ij}\gamma_i(x)\gamma_j(y)
$$
로 쓸 수 있으며, 여기서 $c_{ij} \in \mathbf{Z}$는 어떤
보편 상수들이다. 다시 이 등식이 다항식환 $\mathbf{Q}[x, y]$에서
성립한다는 사실로부터 모든 계수 $c_{ij}$가 원하는 대로
$(nm)!/n!(m!)^n$과 같음을 얻는다.
\end{proof}

\begin{lemma}
\label{dpa-lemma-two-ideals}
$A$를 두 아이디얼 $I, J \subset A$를 갖는 환이라 하자.
$\gamma$를 $I$ 위의 분할 거듭제곱 구조,
$\delta$를 $J$ 위의 분할 거듭제곱 구조라 하자. 그러면
\begin{enumerate}
\item $\gamma$와 $\delta$는 $IJ$ 위에서 일치한다.
\item $\gamma$와 $\delta$가 $I \cap J$ 위에서 일치하면,
이들은 $I + J$ 위의 유일한 분할 거듭제곱 구조 $\epsilon$의
제한이다.
\end{enumerate}
\end{lemma}

\begin{proof}
$x \in I$, $y \in J$라 하자. 그러면
$$
\gamma_n(xy) = y^n\gamma_n(x) = n! \delta_n(y) \gamma_n(x) =
\delta_n(y) x^n = \delta_n(xy).
$$
따라서 $\gamma$와 $\delta$는 $IJ$의 (가법적) 생성자 집합
위에서 일치한다. Definition \ref{dpa-definition-divided-powers}의
성질 (4)에 의해 이들은 $IJ$ 전체에서 일치한다.

\medskip\noindent
$\gamma$와 $\delta$가 $I \cap J$ 위에서 일치한다고 가정하자.
$z \in I + J$라 하고, $x \in I$, $y \in J$를 택하여
$z = x + y$로 쓰자. 그러면 모든 $n \geq 1$에 대해
$$
\epsilon_n(z) = \sum \gamma_i(x)\delta_{n - i}(y)
$$
로 놓는다. 이것이 잘 정의됨을 보기 위해 $x' \in I$,
$y' \in J$이고 $z = x' + y'$인 다른 표현이 주어졌다고 하자.
그러면 $w = x - x' = y' - y \in I \cap J$이다. 따라서 원하는 대로
\begin{align*}
\sum\nolimits_{i + j = n} \gamma_i(x)\delta_j(y)
& =
\sum\nolimits_{i + j = n} \gamma_i(x' + w)\delta_j(y) \\
& =
\sum\nolimits_{i' + l + j = n} \gamma_{i'}(x')\gamma_l(w)\delta_j(y) \\
& =
\sum\nolimits_{i' + l + j = n} \gamma_{i'}(x')\delta_l(w)\delta_j(y) \\
& =
\sum\nolimits_{i' + j' = n} \gamma_{i'}(x')\delta_{j'}(y + w) \\
& =
\sum\nolimits_{i' + j' = n} \gamma_{i'}(x')\delta_{j'}(y')
\end{align*}
이다. 이로써 모든 $n \geq 1$에 대해 사상
$\epsilon_n : I + J \to I + J$를 정의했다.
$\epsilon_n \mid_{I} = \gamma_n$과
$\epsilon_n \mid_{J} = \delta_n$임은 쉽게 알 수 있다.
이제 사상들의 모음 $\epsilon_n$에 대해 Definition
\ref{dpa-definition-divided-powers}의 조건 (1)--(5)를 증명한다.
성질 (1)과 (3)은 자명하다. (4)를 보기 위해
$x, x' \in I$, $y, y' \in J$이고 $z = x + y$,
$z' = x' + y'$라고 하여 계산하면
\begin{align*}
\epsilon_n(z + z') & =
\sum\nolimits_{a + b = n} \gamma_a(x + x')\delta_b(y + y') \\
& =
\sum\nolimits_{i + i' + j + j' = n}
\gamma_i(x) \gamma_{i'}(x')\delta_j(y)\delta_{j'}(y') \\
& =
\sum\nolimits_{k = 0, \ldots, n}
\sum\nolimits_{i+j=k} \gamma_i(x)\delta_j(y)
\sum\nolimits_{i'+j'=n-k} \gamma_{i'}(x')\delta_{j'}(y') \\
& =
\sum\nolimits_{k = 0, \ldots, n}\epsilon_k(z)\epsilon_{n-k}(z')
\end{align*}
을 얻는다. 이제 Lemma \ref{dpa-lemma-check-on-generators}에 의해
$I$ 또는 $J$의 원소들에 대해 (2)와 (5)를 증명하는 것으로
충분하다. $\gamma$와 $\delta$가 분할 거듭제곱 구조이므로
이는 자명하다.

\medskip\noindent
이로써 $I+J$ 위에, $I$와 $J$로 제한하면 각각 $\gamma$와
$\delta$가 되는 분할 거듭제곱 구조 $\epsilon$의 존재가
증명되었다. 그 유일성은 상당히 자명하다.
\end{proof}

\begin{lemma}
\label{dpa-lemma-nil}
$p$를 소수라 하자. $A$를 환, $I \subset A$를 아이디얼,
$\gamma$를 $I$ 위의 분할 거듭제곱 구조라 하자.
$A/I$에서 $p$가 멱영이라고 가정하자. 그러면 $I$가 국소
멱영이기 위한 필요충분조건은 $A$에서 $p$가 멱영인 것이다.
\end{lemma}

\begin{proof}
$A$에서 $p^N = 0$이면, $x \in I$에 대해
$(pN)!$이 $p^N$으로 나누어지므로
$x^{pN} = (pN)!\gamma_{pN}(x) = 0$이다.
거꾸로 $I$가 국소 멱영이라고 가정하자. 또한 $A/I$에서 $p$가
멱영이라고 가정했으므로, 어떤 $r$에 대해 $p^r \in I$이다.
따라서 $p^r$이 멱영이고, 따라서 $p$도 멱영이다.
\end{proof}







\section{분할 거듭제곱 환}
\label{dpa-section-divided-power-rings}

\noindent
분할 거듭제곱 환들의 범주가 있다. 그 정의는 다음과 같다.

\begin{definition}
\label{dpa-definition-divided-power-ring}
{\it 분할 거듭제곱 환}은 삼중항 $(A, I, \gamma)$로서,
$A$는 환이고 $I \subset A$는 아이디얼이며
$\gamma = (\gamma_n)_{n \geq 1}$은 $I$ 위의 분할 거듭제곱
구조인 것이다. {\it 분할 거듭제곱 환의 준동형}
$\varphi : (A, I, \gamma) \to (B, J, \delta)$은
$\varphi(I) \subset J$이고 모든 $x \in I$와 $n \geq 1$에 대해
$\delta_n(\varphi(x)) = \varphi(\gamma_n(x))$인 환 준동형
$\varphi : A \to B$이다.
\end{definition}

\noindent
때때로 “$(B, J, \delta)$를 $(A, I, \gamma)$ 위의 분할
거듭제곱 대수라 하자”라고 말한다. 이는 $(B, J, \delta)$가
분할 거듭제곱 환이고 분할 거듭제곱 환의 준동형
$(A, I, \gamma) \to (B, J, \delta)$이 함께 주어졌음을 뜻한다.

\begin{lemma}
\label{dpa-lemma-limits}
분할 거듭제곱 환의 범주는 모든 극한을 가지며, 이 극한들은
환의 범주에서의 극한들과 일치한다.
\end{lemma}

\begin{proof}
빈 극한은 영환이다(이상하지만 우리에게는 이것이 필요하다).
분할 거듭제곱 환들의 모음 $(A_t, I_t, \gamma_t)$,
$t \in T$의 곱은 $(\prod A_t, \prod I_t, \gamma)$로 주어지며,
여기서 $\gamma_n((x_t)) = (\gamma_{t, n}(x_t))$이다.
$\alpha, \beta : (A, I, \gamma) \to (B, J, \delta)$의 등화자는
아이디얼 $C \cap I$와 유도된 분할 거듭제곱을 갖는
$C = \{a \in A \mid \alpha(a) = \beta(a)\}$이다.
따라서 모든 극한이 존재한다. Categories, Lemma
\ref{categories-lemma-limits-products-equalizers}를 보라.
\end{proof}

\noindent
다음 보조정리는 분할 거듭제곱 대수의 경우에 매우 일반적인
범주론적 현상을 보여 준다.

\begin{lemma}
\label{dpa-lemma-a-version-of-brown}
$\mathcal{C}$를 분할 거듭제곱 환의 범주라 하고
$F : \mathcal{C} \to \textit{Sets}$를 함자라 하자.
다음을 가정하자.
\begin{enumerate}
\item 어떤 기수 $\kappa$가 존재하여, 모든
$f \in F(A, I, \gamma)$에 대해 $\mathcal{C}$의 사상
$(A', I', \gamma') \to (A, I, \gamma)$과
$f' \in F(A', I', \gamma')$가 존재하고, $f$는 이 원소의 상이며
$|A'| \leq \kappa$이다.
\item $F$는 극한과 가환한다.
\end{enumerate}
그러면 $F$는 표현가능하다. 즉 $\mathcal{C}$의 객체
$(B, J, \delta)$가 존재하여 $(A, I, \gamma)$에 관해 함자적으로
$$
F(A, I, \gamma) = \Hom_\mathcal{C}((B, J, \delta), (A, I, \gamma))
$$
이다.
\end{lemma}

\begin{proof}
이는 Categories, Lemma
\ref{categories-lemma-a-version-of-brown}의 특수한 경우이다.
\end{proof}

\begin{lemma}
\label{dpa-lemma-colimits}
분할 거듭제곱 환의 범주는 모든 여극한을 갖는다.
\end{lemma}

\begin{proof}
빈 여극한은 분할 거듭제곱 아이디얼 $(0)$을 갖는
$\mathbf{Z}$이다. 일반적인 여극한을 논하자.
$\mathcal{C}$를 범주라 하고
$c \mapsto (A_c, I_c, \gamma_c)$를 도식이라 하자. 함자
$$
F(B, J, \delta) = \lim_{c \in \mathcal{C}}
Hom((A_c, I_c, \gamma_c), (B, J, \delta))
$$
를 생각하자. 임의의
$f = (f_c)_{c \in C} \in F(B, J, \delta)$는 다음 성질을 갖는다.
모든 상 $f_c(A_c)$는 $B$의 유계 기수 $\kappa$인 부분환
$B'$을 생성하고, 모든 상 $f_c(I_c)$는 $B'$의 분할 거듭제곱
부분아이디얼 $J'$을 생성한다. 그리고 $f$는
$F(B')$의 원소 $f'$ 뒤에 포함 $B' \to B$를 합성하는 것으로
분해된다. 또한 $F$는 극한과 가환한다. 따라서 Lemma
\ref{dpa-lemma-a-version-of-brown}을 적용하면 $F$가 표현가능함을
알 수 있고, 이로써 끝난다.
\end{proof}

\begin{remark}
\label{dpa-remark-pushout}
분할 거듭제곱 환의 범주에서 다음 도식이 밂이라고 하자
(Lemma \ref{dpa-lemma-colimits}에 의해 밂은 존재한다).
$$
\xymatrix{
(B, J, \delta) \ar[r] & (B'', J'', \delta'') \\
(A, I, \gamma) \ar[r] \ar[u] & (B', J', \delta') \ar[u]
}
$$
그러면
\begin{enumerate}
\item $B''/J'' = B/J \otimes_{A/I} B'/J'$이다.
\item 사상 $\varphi : B \otimes_A B' \to B''$는 전사이다.
\item $\varphi$가 유도하는 사상
$J \otimes_A B' + B \otimes_A J' \to J''$는 전사이다.
\end{enumerate}
(1)을 보려면 사상 $(B'', J'', \delta'') \to (C, (0), \emptyset)$을
생각하고 밂의 정의를 적용한다. (2)를 보려면 상
$\tilde B = \varphi(B \otimes_A B')$와 아이디얼
$\tilde J = \varphi(J \otimes_A B' + B \otimes_A J')$을
생각한다. $\delta''$가 $\delta$ 및 $\delta'$와 양립하므로,
모든 $n \geq 1$에 대해
$\delta''_n(\tilde J) \subset \tilde J$임은 자명하다.
따라서 밂의 보편성에 의해 $\tilde B = B''$와
$\tilde J = J''$이어야 한다.
\end{remark}

\begin{remark}
\label{dpa-remark-forgetful}
망각 함자 $(A, I, \gamma) \mapsto A$는 여극한과 가환하지 않는다.
예를 들어 Remark \ref{dpa-remark-pushout}에서 논한 것처럼,
분할 거듭제곱 환의 범주에서 다음 도식이 밂이라고 하자.
$$
\xymatrix{
(B, J, \delta) \ar[r] & (B'', J'', \delta'') \\
(A, I, \gamma) \ar[r] \ar[u] & (B', J', \delta') \ar[u]
}
$$
그러면 일반적으로 사상 $B \otimes_A B' \to B''$은 동형이 아니다.
구체적인 예는
$(A, I, \gamma) = (\mathbf{Z}, (0), \emptyset)$,
$(B, J, \delta) = (\mathbf{Z}/4\mathbf{Z}, 2\mathbf{Z}/4\mathbf{Z}, \delta)$,
그리고
$(B', J', \delta') =
(\mathbf{Z}/4\mathbf{Z}, 2\mathbf{Z}/4\mathbf{Z}, \delta')$로
주어지며, 여기서 $\delta_2(2) = 2$이고
$\delta'_2(2) = 0$이다. 더 정확히 말하면 Lemma
\ref{dpa-lemma-need-only-gamma-p}를 사용하여, $\delta$와
각각 $\delta'$을 $J$와 각각 $J'$ 위의 유일한 분할 거듭제곱
구조로 잡는다. 이때 $\delta_2 : J \to J$와 각각
$\delta'_2 : J' \to J'$는 사상
$0 \mapsto 0, 2 \mapsto 2$와 각각
$0 \mapsto 0, 2 \mapsto 0$이다.
그러면 $(B'', J'', \delta'') = (\mathbf{F}_2, (0), \emptyset)$이고,
이는 텐서곱과 일치하지 않는다.
\end{remark}


\section{분할 거듭제곱의 연장}
\label{dpa-section-extend}

\noindent
정의는 다음과 같다.

\begin{definition}
\label{dpa-definition-extends}
분할 거듭제곱 환 $(A, I, \gamma)$과 환 사상
$A \to B$가 주어졌다고 하자. $IB$ 위에 분할 거듭제곱 구조
$\bar \gamma$가 존재하여
$(A, I, \gamma) \to (B, IB, \bar\gamma)$가 분할 거듭제곱 환의
준동형이 되면, $\gamma$가 $B$로 {\it 연장된다}고 한다.
\end{definition}

\begin{lemma}
\label{dpa-lemma-gamma-extends}
$(A, I, \gamma)$를 분할 거듭제곱 환이라 하고,
$A \to B$를 환 사상이라 하자.
$\gamma$가 $B$로 연장된다면 그 연장은 유일하다.
다음 조건 가운데 적어도 하나가 성립한다고 가정하자.
\begin{enumerate}
\item $IB = 0$,
\item $I$가 주 아이디얼이다. 또는
\item $A \to B$가 평탄하다.
\end{enumerate}
그러면 $\gamma$는 $B$로 연장된다.
\end{lemma}

\begin{proof}
$IB$의 임의의 원소는 $b_i \in B$ 및 $x_i \in I$를 써서 유한합
$\sum\nolimits_{i=1}^t b_ix_i$로 나타낼 수 있다.
$\gamma$가 $IB$ 위의 $\bar\gamma$로 연장된다면
$\bar\gamma_n(x_i) = \gamma_n(x_i)$이다.
따라서 정의 \ref{dpa-definition-divided-powers}의 조건 (3)과 (4)에 의해
$$
\bar\gamma_n(\sum\nolimits_{i=1}^t b_ix_i) =
\sum\nolimits_{n_1 + \ldots + n_t = n}
\prod\nolimits_{i = 1}^t b_i^{n_i}\gamma_{n_i}(x_i)
$$
이다. 그러므로 $\bar\gamma$가 존재한다면 유일하다.

\medskip\noindent
$IB = 0$이면 $\bar\gamma_n(0) = 0$으로 놓으면 된다.
$I = (x)$이면 $\bar\gamma_n(bx) = b^n\gamma_n(x)$로 정의한다.
이는 잘 정의된다. 실제로 $b'x = bx$, 즉 $(b - b')x = 0$이면
\begin{align*}
b^n\gamma_n(x) - (b')^n\gamma_n(x)
& =
(b^n - (b')^n)\gamma_n(x) \\
& =
(b^{n - 1} + \ldots + (b')^{n - 1})(b - b')\gamma_n(x) = 0
\end{align*}
이다. 이는 $\gamma_n(I) \subset I$이므로 $\gamma_n(x)$가 $x$로
나누어지고, 따라서 $b - b'$에 의해 영이 되기 때문이다.
이제 정의 \ref{dpa-definition-divided-powers}의 조건 (1)--(5)를 증명하자.
(1), (2), (3), (5)는 구성으로부터 명백하다.
(4)에 대해서는 $y, z \in IB$, 이를테면 $y = bx$ 및 $z = cx$라 하자.
그러면 $y + z = (b + c)x$이므로
\begin{align*}
\bar\gamma_n(y + z)
& =
(b + c)^n\gamma_n(x) \\
& =
\sum \frac{n!}{i!(n - i)!}b^ic^{n -i}\gamma_n(x) \\
& =
\sum b^ic^{n - i}\gamma_i(x)\gamma_{n - i}(x) \\
& =
\sum \bar\gamma_i(y)\bar\gamma_{n -i}(z)
\end{align*}
를 얻으며, 이는 원하는 바이다.

\medskip\noindent
$A \to B$가 평탄하다고 가정하자.
$b_1, \ldots, b_r \in B$ 및 $x_1, \ldots, x_r \in I$라 하자.
$\bar\gamma_n$이 존재한다면
$$
\bar\gamma_n(\sum b_ix_i) =
\sum b_1^{e_1} \ldots b_r^{e_r} \gamma_{e_1}(x_1) \ldots \gamma_{e_r}(x_r)
$$
이며, 여기서 합은 $e_1 + \ldots + e_r = n$인 경우를 모두 돈다.
이제 $b_i = \sum a_{ij}c_j$를 만족하는 $c_1, \ldots, c_s \in B$와
$a_{ij} \in A$가 있다고 하자.
$y_j = \sum a_{ij}x_i$로 놓으면
$$
\sum b_1^{e_1} \ldots b_r^{e_r} \gamma_{e_1}(x_1) \ldots \gamma_{e_r}(x_r) =
\sum c_1^{d_1} \ldots c_s^{d_s} \gamma_{d_1}(y_1) \ldots \gamma_{d_s}(y_s)
$$
가 $B$에서 성립한다고 주장한다. 여기서 오른쪽 합은
$d_1 + \ldots + d_s = n$인 경우를 모두 돈다.
실제로 분할 거듭제곱 구조의 공리를 사용하면 양변을
$\mathbf{Z}[a_{ij}]$에 속하는 계수와
$c_1^{d_1} \ldots c_s^{d_s}\gamma_{e_1}(x_1) \ldots \gamma_{e_r}(x_r)$
꼴의 항들로 이루어진 합으로 전개할 수 있다.
계수들이 일치함을 보려면, $\gamma$가 $(x_1, \ldots, x_r)$ 위의
유일한 분할 거듭제곱 구조일 때 이 결과가
$\mathbf{Q}[x_1, \ldots, x_r, c_1, \ldots, c_s, a_{ij}]$에서
성립한다는 점에 주목하면 된다.
Lazard 정리(Algebra, Theorem \ref{algebra-theorem-lazard})에 의해
$B$를 유한 자유 $A$-가군들의 유향 여극한으로 나타낼 수 있다.
특히 $z \in IB$가 $z = \sum x_ib_i$ 및
$z = \sum x'_{i'}b'_{i'}$로 나타난다면,
$b_i = \sum a_{ij}c_j$ 및 $b'_{i'} = \sum a'_{i'j}c_j$를 만족하고
$y_j = \sum x_ia_{ij} = \sum x'_{i'}a'_{i'j}$가 성립하도록 하는
$c_1, \ldots, c_s \in B$와 $a_{ij}, a'_{i'j} \in A$를 찾을 수 있다
\footnote{이는 Algebra, Theorem \ref{algebra-theorem-lazard}을 사용하지
않고도 증명할 수 있다. 실제로 $z = \sum x_ib_i$이고
$z = \sum x'_{i'}b'_{i'}$이면
$\sum x_ib_i - \sum x'_{i'}b'_{i'} = 0$은 $A$-가군 $B$에서의 관계이다.
따라서 Algebra, Lemma \ref{algebra-lemma-flat-eq}를 적용하되,
$f_i$ 대신 $x_i$와 $x'_{i'}$를, 그 안의 $x_i$ 역할에는
$b_i$와 $b'_{i'}$를 사용하면, 필요한 $c_1, \ldots, c_s \in B$와
$a_{ij}, a'_{i'j} \in A$의 존재를 얻는다.}.
그러므로 위 절차는 $IB$ 위의 잘 정의된 사상 $\bar\gamma_n$을 준다.
구성에 의해 $\bar\gamma$는 조건 (1), (3), (4)를 만족한다.
더욱이 $x \in I$에 대해 $\bar\gamma_n(x) = \gamma_n(x)$이다.
따라서 보조정리 \ref{dpa-lemma-check-on-generators}에 의해 $\bar\gamma$는
$IB$ 위의 분할 거듭제곱 구조이다.
\end{proof}

\begin{lemma}
\label{dpa-lemma-kernel}
$(A, I, \gamma)$를 분할 거듭제곱 환이라 하자.
\begin{enumerate}
\item $\varphi : (A, I, \gamma) \to (B, J, \delta)$가 분할 거듭제곱 환의
준동형이면, 모든 $n \geq 1$에 대해 $\Ker(\varphi) \cap I$는
$\gamma_n$에 의해 보존된다.
\item $\mathfrak a \subset A$를 아이디얼이라 하고
$I' = I \cap \mathfrak a$로 놓자. 다음 조건들은 동치이다.
\begin{enumerate}
\item 모든 $n > 0$에 대해 $I'$는 $\gamma_n$에 의해 보존된다.
\item $\gamma$는 $A/\mathfrak a$로 연장된다.
\item 아이디얼로서 $I'$의 생성원 집합 $x_i$가 존재하여,
모든 $n > 0$에 대해 $\gamma_n(x_i) \in I'$이다.
\end{enumerate}
\end{enumerate}
\end{lemma}

\begin{proof}
(1)의 증명. 이는 명백하다. (2)(a)를 가정하자. $x \in I$에 대해
$\bar\gamma_n(x \bmod I') = \gamma_n(x) \bmod I'$로 정의한다.
가정에 의해 $\gamma_j(y) \in I'$이고 정의
\ref{dpa-definition-divided-powers} (4)가 성립하므로, $y \in I'$에 대해
$\gamma_n(x + y) = \gamma_n(x) \bmod I'$이다. 따라서 이는 잘 정의된다.
$\gamma$가 분할 거듭제곱 구조이므로 $\bar\gamma$도 그러한 구조임이 명백하다.
따라서 (2)(b)가 성립한다. 또한 (1)에 의해 (2)(b)는 (2)(a)를 함의한다.
(2)(a)가 (2)(c)를 함의함은 명백하다. (2)(c)를 가정하자.
$x = ax_i$일 때 $\gamma_n(x) = a^n\gamma_n(x_i) \in I'$임에 주목하자.
따라서 아벨군으로서 $I'$의 한 생성원 집합에 대해
$\gamma_n(x) \in I'$임을 알 수 있다.
이 생성원들로 된 표현의 길이에 대한 귀납법에 의해,
$\forall n : \gamma_n(x), \gamma_n(y) \in I'$이면
$\forall n : \gamma_n(x + y) \in I'$임을 보이는 것으로 충분하다.
이는 분할 거듭제곱 구조의 네 번째 공리에서 곧바로 따른다.
\end{proof}

\begin{lemma}
\label{dpa-lemma-sub-dp-ideal}
$(A, I, \gamma)$를 분할 거듭제곱 환이라 하고,
$E \subset I$를 부분집합이라 하자.
$\gamma$에 의해 보존되고 모든 $f \in E$를 포함하는 가장 작은
아이디얼 $J \subset I$는 $\gamma_n(f)$, $n \geq 1$, $f \in E$로
생성되는 아이디얼 $J$이다.
\end{lemma}

\begin{proof}
보조정리 \ref{dpa-lemma-kernel}에서 곧바로 따른다.
\end{proof}

\begin{lemma}
\label{dpa-lemma-extend-to-completion}
$(A, I, \gamma)$를 분할 거듭제곱 환이라 하고 $p$를 소수라 하자.
$p$가 $A/I$에서 멱영이면 다음이 성립한다.
\begin{enumerate}
\item $p$-진 완비화 $A^\wedge = \lim_e A/p^eA$는 $A/I$로 전사한다.
\item 이 사상의 핵은 $I$의 $p$-진 완비화 $I^\wedge$이다.
\item 각 $\gamma_n$은 $p$-진 위상에 관해 연속이고
$\gamma_n^\wedge : I^\wedge \to I^\wedge$로 연장되어,
$I^\wedge$ 위의 분할 거듭제곱 구조를 정의한다.
\end{enumerate}
더 나아가 $A$가 $\mathbf{Z}_{(p)}$-대수이면 다음도 성립한다.
\begin{enumerate}
\item[(4)] 충분히 큰 $e$에 대해 아이디얼 $p^eA \subset I$는
분할 거듭제곱 구조 $\gamma$에 의해 보존되고,
$$
(A^\wedge, I^\wedge, \gamma^\wedge) = \lim_e (A/p^eA, I/p^eA, \bar\gamma)
$$
가 분할 거듭제곱 환의 범주에서 성립한다.
\end{enumerate}
\end{lemma}

\begin{proof}
$p^tA/I = 0$, 즉 $p^tA \subset I$가 되도록 하는 정수
$t \geq 1$을 택하자. 사상 $A^\wedge \to A/I$는 합성
$A^\wedge \to A/p^tA \to A/I$이며, 이는 전사이다(예를 들어
Algebra, Lemma \ref{algebra-lemma-completion-generalities}에 의한다).
$e \geq t$에 대해 $p^eI \subset p^eA \cap I \subset p^{e - t}I$이므로,
합성 $A^\wedge \to A/I$의 핵은 $I$의 $p$-진 완비화이다.
사상 $\gamma_n$은 다음 식과 분할 거듭제곱 구조의 공리에 의해 연속이다.
$$
\gamma_n(x + p^ey) =
\sum\nolimits_{i + j = n} p^{je}\gamma_i(x)\gamma_j(y) =
\gamma_n(x) \bmod p^eI
$$
사상 $\gamma_n^\wedge$들이 $\gamma_n$들로부터 분할 거듭제곱 구조의
공리들을 물려받는다는 것은 명백하다. 마지막 주장을 보이기 위해
$e > t$라 하자. 그러면 $p^eA \subset I$이고
$\gamma_1(p^eA) \subset p^eA$이며, $n > 1$에 대해
$$
\gamma_n(p^ea) = p^n \gamma_n(p^{e - 1}a) = \frac{p^n}{n!} p^{n(e - 1)}a^n
\in p^e A
$$
이다. 이는 $p^n/n! \in \mathbf{Z}_{(p)}$이고 $n \geq 2$ 및
$e \geq 2$이므로 $n(e - 1) \geq e$이기 때문이다.
이로써 $\gamma$가 $A/p^eA$로 연장됨을 증명했으며, 보조정리
\ref{dpa-lemma-kernel}을 보라. 극한에 관한 주장은 보조정리
\ref{dpa-lemma-limits}의 증명에 나오는 극한의 구성으로부터 명백하다.
\end{proof}




\section{분할 거듭제곱 다항식 대수}
\label{dpa-section-divided-power-polynomial-ring}

\noindent
매우 유용한 예로 {\it 분할 거듭제곱 다항식 대수}가 있다.
$A$를 환이라 하고 $t \geq 1$이라 하자. 다음 $A$-대수를
$A\langle x_1, \ldots, x_t \rangle$로 나타내기로 한다.
$A$-가군으로서는
$$
A\langle x_1, \ldots, x_t \rangle =
\bigoplus\nolimits_{n_1, \ldots, n_t \geq 0} A x_1^{[n_1]} \ldots x_t^{[n_t]}
$$
로 놓고, 곱셈은
$$
x_i^{[n]}x_i^{[m]} = \frac{(n + m)!}{n!m!}x_i^{[n + m]}.
$$
로 주어진다. 또한 $x_i = x_i^{[1]}$로 놓는다.
$1 = x_1^{[0]} \ldots x_t^{[0]}$임에 유의하라.
비슷한 구성으로 무한히 많은 변수의 분할 거듭제곱 다항식 대수를
얻을 수 있다. $n > 0$일 때 $x_i^{[n]}$을 영으로 보내는 표준
$A$-대수 사상
$A\langle x_1, \ldots, x_t \rangle \to A$가 있다.
이 사상의 핵을 $A\langle x_1, \ldots, x_t \rangle_{+}$로 나타낸다.

\begin{lemma}
\label{dpa-lemma-divided-power-polynomial-algebra}
$(A, I, \gamma)$를 분할 거듭제곱 환이라 하자. 다음 아이디얼 위에는
유일한 분할 거듭제곱 구조 $\delta$가 존재한다.
$$
J = IA\langle x_1, \ldots, x_t \rangle + A\langle x_1, \ldots, x_t \rangle_{+}
$$
이 구조는 다음을 만족한다.
\begin{enumerate}
\item $\delta_n(x_i) = x_i^{[n]}$.
\item $(A, I, \gamma) \to (A\langle x_1, \ldots, x_t \rangle, J, \delta)$
는 분할 거듭제곱 환의 준동형이다.
\end{enumerate}
더욱이 $(A\langle x_1, \ldots, x_t \rangle, J, \delta)$는 다음
보편 성질을 갖는다. 분할 거듭제곱 환의 준동형
$\varphi : (A\langle x_1, \ldots, x_t \rangle, J, \delta) \to
(C, K, \epsilon)$를 주는 것은 분할 거듭제곱 환의 준동형
$A \to C$와 원소 $k_1, \ldots, k_t \in K$를 주는 것과 같다.
\end{lemma}

\begin{proof}
변수가 하나인 분할 거듭제곱 다항식 대수의 경우에 보조정리를
증명하겠다. 일반적인 경우도 정확히 같은 방법으로 논증할 수 있으며,
또는 $A\langle x_1, \ldots, x_t\rangle$가 분할 거듭제곱 변수
$x_1, \ldots, x_t$를 하나씩 첨가하여 얻은 환과 동형임을 이용할 수 있다.

\medskip\noindent
$A\langle x \rangle_{+}$를 $x, x^{[2]}, x^{[3]}, \ldots$로
생성되는 아이디얼이라 하자.
$J = IA\langle x \rangle + A\langle x \rangle_{+}$이고
$$
IA\langle x \rangle \cap A\langle x \rangle_{+} =
IA\langle x \rangle \cdot A\langle x \rangle_{+}
$$
임에 유의하라. 따라서 보조정리 \ref{dpa-lemma-two-ideals}에 의해,
아이디얼 $IA\langle x \rangle$과 $A\langle x \rangle_{+}$ 위에
분할 거듭제곱 구조가 존재함을 보이면 충분하다.
$A \to A\langle x \rangle$가 평탄하므로 첫 번째 구조의 존재는
보조정리 \ref{dpa-lemma-gamma-extends}에서 따른다.
두 번째 구조에 대해서는 $A$가 비틀림이 없으면 보조정리
\ref{dpa-lemma-silly} (4)를 적용하여 $\delta$가 존재함을 알 수 있다.
실제로 원소 $x^{[m]}$들을 생성원으로 택하면
$(x^{[m]})^n = \frac{(nm)!}{(m!)^n} x^{[nm]}$
이고, $n!$은 정수 $\frac{(nm)!}{(m!)^n}$을 나눈다.
일반적으로 $A = R/\mathfrak a$라 쓰자. 여기서 $R$은 비틀림이 없는
어떤 환이다(예를 들어 $\mathbf{Z}$ 위의 다항식 환).
$R\langle x \rangle \to A\langle x \rangle$의 핵은
$\bigoplus \mathfrak a x^{[m]}$이다. 보조정리 \ref{dpa-lemma-kernel}의
판정법 (2)(c)를 적용하면 $R\langle x \rangle_{+}$ 위의
분할 거듭제곱 구조가 원하는 대로 $A\langle x \rangle$로 연장됨을 알 수 있다.

\medskip\noindent
보편 성질의 증명. 분할 거듭제곱 환의 준동형
$\varphi : A \to C$와 $k_1, \ldots, k_t \in K$가 주어지면
$$
A\langle x_1, \ldots, x_t \rangle \to C,\quad
x_1^{[n_1]} \ldots x_t^{[n_t]} \longmapsto
\epsilon_{n_1}(k_1) \ldots \epsilon_{n_t}(k_t)
$$
를 생각하되 계수에는 $\varphi$를 사용한다.
확인할 것은 이것이 $A$-대수 준동형이라는 사실뿐이다
(세부 사항은 생략한다). 역구성은 명백하다.
\end{proof}

\begin{remark}
\label{dpa-remark-divided-power-polynomial-algebra}
$(A, I, \gamma)$를 분할 거듭제곱 환이라 하자.
무한히 많은 변수에 대한 보조정리
\ref{dpa-lemma-divided-power-polynomial-algebra}의 변형이 있다.
먼저 $s < t$이면 표준 사상
$$
A\langle x_1, \ldots, x_s \rangle \to A\langle x_1, \ldots, x_t\rangle
$$
이 있음에 유의하라. 따라서 $W$가 임의의 집합이면
$$
A\langle x_w: w \in W\rangle =
\colim_{E \subset W} A\langle x_e:e \in E\rangle
$$
로 놓는다(여극한은 $W$의 유한 부분집합 $E$들을 돌며,
전이 사상들은 위와 같다). 여극한의 정의에 의해
$A\langle x_w: w \in W\rangle$의 보편 사상 성질은 보조정리
\ref{dpa-lemma-divided-power-polynomial-algebra}에서 서술한 사상 성질과
완전히 유사하다.
\end{remark}

\noindent
다음 보조정리는 \cite{BO}에서 찾을 수 있다.

\begin{lemma}
\label{dpa-lemma-need-only-gamma-p}
$p$를 소수라 하자. $p$로 나누어지지 않는 모든 정수 $n$이 가역인
환을 $A$라 하자. 즉, $A$는 $\mathbf{Z}_{(p)}$-대수이다.
$I \subset A$를 아이디얼이라 하자. $I$ 위의 두 분할 거듭제곱 구조
$\gamma, \gamma'$가 같은 것은 $\gamma_p = \gamma'_p$인 것과 동치이다.
더욱이 사상 $\delta : I \to I$가 주어져 다음을 만족한다고 하자.
\begin{enumerate}
\item 모든 $x \in I$에 대해 $p!\delta(x) = x^p$이다\footnote{이 조건은
나머지 두 조건에서 따른다. Remark \ref{dpa-remark-ryo-suzuki}를 보라.}.
\item 모든 $a \in A$, $x \in I$에 대해 $\delta(ax) = a^p\delta(x)$이다.
\item
$\delta(x + y) =
\delta(x) +
\sum\nolimits_{i + j = p, i,j \geq 1} \frac{1}{i!j!} x^i y^j +
\delta(y)$가 모든 $x, y \in I$에 대해 성립한다.
\end{enumerate}
그러면 $\gamma_p = \delta$를 만족하는 $I$ 위의 유일한
분할 거듭제곱 구조 $\gamma$가 존재한다.
\end{lemma}

\begin{proof}
$n$이 $p$로 나누어지지 않으면
$\gamma_n(x) = c x \gamma_{n - 1}(x)$이며,
여기서 $c$는 $\mathbf{Z}_{(p)}$의 단위원이다. 또한
$$
\gamma_{pm}(x) = c \gamma_m(\gamma_p(x))
$$
이며, 여기서도 $c$는 $\mathbf{Z}_{(p)}$의 단위원이다.
따라서 첫 번째 주장은 명백하다.
두 번째 주장을 위해서는 이 등식들을 거꾸로 사용하여 모든
$\gamma_n$을 정의할 수 있다. 더 정확히 말해
$n = a_0 + a_1p + \ldots + a_e p^e$이고
$a_i \in \{0, \ldots, p - 1\}$이면
$$
\gamma_n(x) = c_n x^{a_0} \delta(x)^{a_1} \ldots \delta^e(x)^{a_e}
$$
로 놓는다. 여기서 $c_n \in \mathbf{Z}_{(p)}$는
$$
c_n =
{(p!)^{a_1 + a_2(1 + p) + \ldots + a_e(1 + \ldots + p^{e - 1})}}/{n!}.
$$
로 정의된다. 이제 분할 거듭제곱 구조의 공리 (1)--(5)를 보여야 한다.
정의 \ref{dpa-definition-divided-powers}를 보라. (1)과 (3)은 즉시 성립한다.
(2)와 (5)는 직접 계산으로 확인할 수 있으며, 그 계산은 생략한다.
$x, y \in I$라 하자. 다음과 같은 환 사상이 있다고 주장한다.
$$
\varphi : \mathbf{Z}_{(p)}\langle u, v \rangle \longrightarrow A
$$
이 사상은 $u^{[n]}$을 $\gamma_n(x)$로,
$v^{[n]}$을 $\gamma_n(y)$로 보낸다.
$\mathbf{Z}_{(p)}\langle u, v \rangle$의 구성에 따라 확인해야 할 것은
$$
\gamma_n(x)\gamma_m(x) = \frac{(n + m)!}{n!m!} \gamma_{n + m}(x)
$$
가 $A$에서 성립하는지와 $y$에 대한 같은 등식이다.
이는 $\gamma$에 대해 (2)가 성립하므로 참이다.
$\epsilon$을 $\mathbf{Z}_{(p)}\langle u, v\rangle$의 아이디얼
$\mathbf{Z}_{(p)}\langle u, v\rangle_{+}$ 위의
분할 거듭제곱 구조라 하자.
다음으로 모든 $n$ 및
$f \in \mathbf{Z}_{(p)}\langle u, v\rangle_{+}$에 대해
$\varphi(\epsilon_n(f)) = \gamma_n(\varphi(f))$라고 주장한다.
$n = 0, 1, \ldots, p - 1$인 경우에는 명백하다.
$n = p$인 경우에는 아이디얼
$\mathbf{Z}_{(p)}\langle u, v\rangle_{+}$의 생성원 집합에 대해
이를 증명하면 충분하다. 이는 $\epsilon_p$와 $\gamma_p = \delta$가
모두 보조정리의 성질 (1), (3)을 만족하기 때문이다.
따라서
$\gamma_p(\gamma_n(x)) = \frac{(pn)!}{p!(n!)^p}\gamma_{pn}(x)$와
$y$에 대한 같은 등식을 증명하면 충분하며,
이는 $\gamma$에 대해 (5)가 성립하므로 따른다.
이제 $n = a_0 + a_1p + \ldots + a_e p^e$가 위와 같은
$p$-진 전개로 쓰인 임의의 정수이면
$$
\epsilon_n(f) =
c_n f^{a_0} \gamma_p(f)^{a_1} \ldots \gamma_p^e(f)^{a_e}
$$
이다. 이는 $\epsilon$이 분할 거듭제곱 구조이기 때문이다.
따라서 모든 $n$에 대해
$\varphi(\epsilon_n(f)) = \gamma_n(\varphi(f))$임을 알 수 있다.
이를 $f = u + v$에 적용하면, $\epsilon$이 분할 거듭제곱 구조라는
사실로부터 $\gamma$에 대한 공리 (4)가 따름을 알 수 있다.
\end{proof}

\begin{remark}
\label{dpa-remark-ryo-suzuki}
보조정리 \ref{dpa-lemma-need-only-gamma-p}에서처럼 $A \supset I$라 하고,
$\delta : I \to I$를 그 보조정리의 (2), (3)을 만족하는 사상이라 하자.
그러면 (1)도 성립한다.
먼저 $n$에 대한 귀납법으로 모든 $n \geq 1$에 대해
$\delta(nx) = n\delta(x) + \frac{n^p - n}{p!}x^p$임을 증명하자.
$n = 1$인 경우에는 성립한다. 어떤 $n$에 대해 성립한다고 가정하면
$\delta((n + 1)x) = \delta(nx) + \delta(x) +
(\sum_{i = 1}^{p - 1} \frac{n^i}{i!(p - i)!})x^p
= (n + 1)\delta(x) + \left( \frac{n^p - n}{p!} + \frac{(n + 1)^p}{p!}
- \frac{n^p + 1}{p!} \right)x^p =
(n + 1)\delta(x) + \frac{(n + 1)^p - (n + 1)}{p!}x^p$,
이므로 $n + 1$에 대해서도 성립한다. 한편
$\delta(nx) = n^p\delta(x)$이므로
$(n^p - n)\delta(x) = \frac{n^p - n}{p!}x^p$이다.
그런데 임의의 $p$에 대해 $p^2 \not \mid n^p - n$을 만족하는
$n$이 존재한다. 실제로
$(\mathbf{Z}/p^2\mathbf{Z})^\times \cong \mathbf{Z}/p(p-1)\mathbf{Z}$.
따라서 $n^{p-1} \not \equiv 1 \mod p^2$인 $n$이 존재한다.
그러므로 $p!\delta(x)=x^p$을 얻는다.
\end{remark}









\section{Tate 분해}
\label{dpa-section-tate}

\noindent
이 절에서는 분할 거듭제곱 구조와 미분 등급 대수를 결합하여
\cite{Tate-homology}와 \cite{AH}에서 구성한 분해를 간략히 논한다.
이 절에서는 미분 등급 대수에 대해 {\it 호몰로지 표기}를 사용한다.
우리의 미분 등급 대수는 음이 아닌 호몰로지 차수에 놓인다.
따라서 미분 등급 대수 $(A, \text{d})$는 곱셈을 갖춘 사슬 복합체
$$
\ldots \to A_2 \to A_1 \to A_0 \to 0 \to \ldots
$$
로 주어진다.

\medskip\noindent
$R$을 환이라 하자(평소와 같이 가환환이다).
이 절에서는 음의 차수 성분이 영인 등급 $R$-대수
$A = \bigoplus_{d \geq 0} A_d$를 자주 다룬다.
$A_+ = \bigoplus_{d > 0} A_d$로 놓고,
$A_{even} = \bigoplus_{d \geq 0} A_{2d}$ 및
$A_{odd} = \bigoplus_{d \geq 0} A_{2d + 1}$로 쓴다.
동차 원소 $x, y$에 대해
$x y = (-1)^{\deg(x)\deg(y)} y x$이면 $A$가 등급 가환임을 상기하자.
더 나아가 홀수 차수의 동차 원소 $x$에 대해 $x^2 = 0$이면
$A$는 엄격히 등급 가환이다. 마지막으로 다음 정의를 이해하기 위해,
$A$가 $\mathbf{Q}$-대수이면 $\gamma_n(x) = x^n/n!$임을 염두에 두자.

\begin{definition}
\label{dpa-definition-divided-powers-graded}
$R$을 환이라 하고, $A = \bigoplus_{d \geq 0} A_d$를 엄격히 등급 가환인
등급 $R$-대수라 하자. 모든 $n > 0$에 대해 정의된 사상들의 모음
$\gamma_n : A_{even, +} \to A_{even, +}$가 다음을 만족하면,
이를 $A$ 위의 {\it 분할 거듭제곱 구조}라 한다.
\begin{enumerate}
\item $x \in A_{2d}$이면 $\gamma_n(x) \in A_{2nd}$이다.
\item 임의의 $x$에 대해 $\gamma_1(x) = x$이고, 또한 $\gamma_0(x) = 1$로 놓는다.
\item $\gamma_n(x)\gamma_m(x) = \frac{(n + m)!}{n! m!} \gamma_{n + m}(x)$,
\item 모든 $x \in A_{even}$ 및 $y \in A_{even, +}$에 대해
$\gamma_n(xy) = x^n \gamma_n(y)$이다.
\item $x, y \in A_{odd}$가 동차이고 $n > 1$이면 $\gamma_n(xy) = 0$이다.
\item $x, y \in A_{even, +}$이면
$\gamma_n(x + y) = \sum_{i = 0, \ldots, n} \gamma_i(x)\gamma_{n - i}(y)$,
\item $x \in A_{even, +}$에 대해 $\gamma_n(\gamma_m(x)) =
\frac{(nm)!}{n! (m!)^n} \gamma_{nm}(x)$이다.
\end{enumerate}
\end{definition}

\noindent
조건 (2), (3), (4), (6), (7)에 의해 $\gamma$는 (가환)환
$A_{even}$의 아이디얼 $A_{even, +}$ 위의 통상적인 분할 거듭제곱
구조임에 유의하라. Sections \ref{dpa-section-divided-powers},
\ref{dpa-section-divided-power-rings},
\ref{dpa-section-extend}, and
\ref{dpa-section-divided-power-polynomial-ring}을 보라.
특히 모든 $x \in A_{even, +}$에 대해 $n! \gamma_n(x) = x^n$이다.
조건 (1)은 $\gamma$가 등급과 양립함을 말하고, 조건 (5)는
$n > 1$일 때 $\gamma_n$이 홀수 차수 동차 원소들의 곱에서
영이 됨을 말한다. 그러나 홀수 차수의 동차 원소
$z_1, z_2, z_3, z_4$에 대해
$$
\gamma_2(z_1 z_2 + z_3 z_4) = z_1z_2z_3z_4
$$
가 영이 아닐 수도 있음에 유의하라.

\begin{example}[홀수 변수의 첨가]
\label{dpa-example-adjoining-odd}
$R$을 환이라 하자. $(A, \gamma)$를 위 정의와 같은 분할 거듭제곱
구조를 갖춘 엄격히 등급 가환인 등급 $R$-대수라 하자.
$d > 0$을 홀수라 하자. 이 상황에서는 차수가 $d$인 변수 $T$를
$A$에 첨가할 수 있다. 즉,
$$
A\langle T \rangle = A \oplus AT
$$
로 놓고, 등급은 $A\langle T \rangle_m = A_m \oplus A_{m - d}T$로 준다.
$A\langle T \rangle$ 위에는 $A$의 주어진 분할 거듭제곱 구조와
양립하는 유일한 분할 거듭제곱 구조가 있다고 주장한다. 즉,
$$
\gamma_n(x + yT) = \gamma_n(x) + \gamma_{n - 1}(x)yT
$$
로 놓는다. 여기서 $x \in A_{even, +}$이고 $y \in A_{odd}$이다.
\end{example}

\begin{example}[짝수 변수의 첨가]
\label{dpa-example-adjoining-even}
$R$을 환이라 하자. $(A, \gamma)$를 위 정의와 같은 분할 거듭제곱
구조를 갖춘 엄격히 등급 가환인 등급 $R$-대수라 하자.
$d > 0$을 짝수라 하자. 이 상황에서는 차수가 $d$인 변수 $T$를
$A$에 첨가할 수 있다. 즉,
$$
A\langle T \rangle = A \oplus AT \oplus AT^{(2)} \oplus AT^{(3)} \oplus \ldots
$$
로 놓고, 곱셈은
$$
T^{(n)} T^{(m)} = \frac{(n + m)!}{n!m!} T^{(n + m)}
$$
로, 등급은
$$
A\langle T \rangle_m =
A_m \oplus A_{m - d}T \oplus A_{m - 2d}T^{(2)} \oplus \ldots
$$
로 준다. $A\langle T \rangle$ 위에는 $A$의 주어진 분할 거듭제곱
구조와 양립하고 $\gamma_n(T^{(i)}) = T^{(ni)}$를 만족하는
유일한 분할 거듭제곱 구조가 있다고 주장한다.
이를 정의하기 위해 먼저
$$
\gamma_n\left(\sum\nolimits_{i > 0} x_i T^{(i)}\right) =
\sum \prod\nolimits_{n = \sum e_i} x_i^{e_i} T^{(ie_i)}
$$
로 놓는다. 여기서 $x_i$는 $A_{even}$에 속한다.
$x_0 \in A_{even, +}$이면
$$
\gamma_n\left(\sum\nolimits_{i \geq 0} x_i T^{(i)}\right) =
\sum\nolimits_{a + b = n}
\gamma_a(x_0)\gamma_b\left(\sum\nolimits_{i > 0} x_iT^{(i)}\right)
$$
로 놓는다. 여기서 $\gamma_b$는 위에서 정의한 바와 같다.
\end{example}

\begin{remark}
\label{dpa-remark-adjoining-set-of-variables}
Examples \ref{dpa-example-adjoining-odd}와
\ref{dpa-example-adjoining-even}에서처럼 하나의 생성원만 첨가하는 대신,
주어진 차수 $d > 0$의 외대수 생성원 또는 분할 거듭제곱 생성원으로
이루어진 집합(무한일 수도 있다)을 첨가할 수도 있다.
구체적으로 Remark \ref{dpa-remark-divided-power-polynomial-algebra}를 따라,
위와 같은 $(A,\gamma)$와 집합 $J$에 대해 $A\langle
T_j:j\in J\rangle$를 $J$의 모든 유한 부분집합 $S$에 걸친 대수
$A\langle T_j:j\in S\rangle$들의 유향 여극한으로 놓는다.
이 대수에는 $A$의 주어진 구조 및 각 생성원 $T_j$ 위의 구조와
양립하는 유일한 분할 거듭제곱 구조가 있음이 즉시 따른다.
\end{remark}

\noindent
이제 분할 거듭제곱 구조의 정의를 미분과 연결한다.
정의를 이해하기 위해 $A$가 $\mathbf{Q}$-대수이고
$x \in A_{even, +}$이면
$\text{d}(x^n/n!) = \text{d}(x) x^{n - 1}/(n - 1)!$임에 유의하라.

\begin{definition}
\label{dpa-definition-divided-powers-dga}
$R$을 환이라 하자. $A = \bigoplus_{d \geq 0} A_d$를 엄격히 등급
가환인 미분 등급 $R$-대수라 하자. $A$ 위의 분할 거듭제곱 구조
$\gamma$가 모든 $x \in A_{even, +}$에 대해
$\text{d}(\gamma_n(x)) = \text{d}(x) \gamma_{n - 1}(x)$를 만족하면,
이 구조는 {\it 미분 등급 구조와 양립한다}고 한다.
\end{definition}

\noindent
주의: $(A, \text{d}, \gamma)$가 정의
\ref{dpa-definition-divided-powers-dga}에서와 같다고 하자.
$x$가 경계이더라도 $\gamma_n(x)$가 경계라는 보장은 없다.
따라서 일반적으로 $\gamma$는 호몰로지 대수 $H(A)$ 위에
분할 거듭제곱 구조를 유도하지 않는다.
일부 논문에서는 이를 보장하기 위해 양립 조건을 하나 더 두지만,
여기서는 그렇게 하지 않기로 한다.

\begin{lemma}
\label{dpa-lemma-dpdga-good}
$(A, \text{d}, \gamma)$와 $(B, \text{d}, \gamma)$가 정의
\ref{dpa-definition-divided-powers-dga}에서와 같다고 하자.
$f : A \to B$를 분할 거듭제곱 구조와 양립하는 미분 등급 대수의
사상이라 하자. 다음을 가정하자.
\begin{enumerate}
\item $k > 0$에 대해 $H_k(A) = 0$이다.
\item $f$는 전사이다.
\end{enumerate}
그러면 $\gamma$는 등급 $R$-대수 $H(B)$ 위에 분할 거듭제곱 구조를 유도한다.
\end{lemma}

\begin{proof}
$x$와 $x'$가 같은 차수 $2d$의 동차 원소이고 $H(B)$에서 같은
호몰로지류를 정의한다고 하자. $x' - x = \text{d}(w)$라 쓰자.
$x$의 올림 $y \in A_{2d}$와 $w$의 올림 $z \in A_{2d + 1}$을 택하자.
그러면 $y' = y + \text{d}(z)$는 $x'$의 올림이다. 따라서
$$
\gamma_n(y') = \sum \gamma_i(y) \gamma_{n - i}(\text{d}(z))
= \gamma_n(y) +
\sum\nolimits_{i < n} \gamma_i(y) \gamma_{n - i}(\text{d}(z))
$$
$A$가 양의 차수에서 비순환이고 모든 $j$에 대해
$\text{d}(\gamma_j(\text{d}(z))) = 0$이므로, 이를
$$
\gamma_n(y') = \gamma_n(y) +
\sum\nolimits_{i < n} \gamma_i(y) \text{d}(z_i)
$$
와 같이 쓸 수 있다. 여기서 $z_i$는 $A$의 원소이다.
더욱이 $0 < i < n$에 대해
$$
\text{d}(\gamma_i(y) z_i) =
\text{d}(\gamma_i(y))z_i + \gamma_i(y)\text{d}(z_i) =
\text{d}(y) \gamma_{i - 1}(y) z_i + \gamma_i(y)\text{d}(z_i)
$$
이다. $\text{d}(y)$가 $B$에서 영으로 가므로 첫 번째 항도 $B$에서
영으로 간다. 따라서 $\gamma_n(x')$와 $\gamma_n(x)$는 $H(B)$의
같은 원소로 간다. 이로써 모든 $d > 0$ 및 $n > 0$에 대해 잘 정의된 사상
$\gamma_n : H_{2d}(B) \to H_{2nd}(B)$를 얻는다.
이것이 $H(B)$ 위의 분할 거듭제곱 구조를 정의한다는 확인은 생략한다.
\end{proof}

\begin{lemma}
\label{dpa-lemma-base-change-div}
$(A, \text{d}, \gamma)$가 정의
\ref{dpa-definition-divided-powers-dga}에서와 같다고 하자.
$R \to R'$를 환 사상이라 하자. 그러면 $\text{d}$와 $\gamma$는
$A' = A \otimes_R R'$ 위에 같은 종류의 구조를 유도하며,
$(A', \text{d}, \gamma)$는 정의
\ref{dpa-definition-divided-powers-dga}에서와 같은 대상이 된다.
\end{lemma}

\begin{proof}
$A'_{even} = A_{even} \otimes_R R'$이고
$A'_{even, +} = A_{even, +} \otimes_R R'$임에 유의하라.
따라서 보일 것은 분할 거듭제곱 $\gamma$가 $A'_{even}$으로
연장된다는 것이다(용어는 정의 \ref{dpa-definition-extends}와 같다).
$\gamma$가 연장됨을 보이고 나면, 이 연장이 원하는 모든 성질을
갖는다는 것은 쉽게 따른다.

\medskip\noindent
임의의 생성원 집합을 갖는 다항식 $R$-대수 $P$와
전사인 $R$-대수 사상 $P \to R'$를 택하자.
환 사상 $A_{even} \to A_{even} \otimes_R P$는 평탄하므로,
보조정리 \ref{dpa-lemma-gamma-extends}에 의해 분할 거듭제곱 $\gamma$는
$A_{even} \otimes_R P$로 유일하게 연장된다.
$J = \Ker(P \to R')$라 하자. $\gamma$가 $A \otimes_R R'$로
연장됨을 보이려면
$I' = \Ker(A_{even, +} \otimes_R P \to A_{even, +} \otimes_R R')$
가 모든 $n > 0$에 대해 $\gamma_n(z) \in I'$인 원소 $z$들로
생성됨을 보이면 충분하다. 이는 $I'$가
$x \in A_{even, +}$ 및 $f \in \Ker(P \to R')$에 대해
$x \otimes f$ 꼴의 원소들로 생성되므로 명백하다.
\end{proof}

\begin{lemma}
\label{dpa-lemma-extend-differential}
$(A, \text{d}, \gamma)$가 정의
\ref{dpa-definition-divided-powers-dga}에서와 같다고 하자.
$d \geq 1$을 정수라 하자. $A\langle T \rangle$를 Example
\ref{dpa-example-adjoining-odd} 또는 \ref{dpa-example-adjoining-even}에서
구성한, $\deg(T) = d$인 $T$에 관한 등급 분할 거듭제곱
다항식 대수라 하자. $f \in A_{d - 1}$가 $\text{d}(f) = 0$을
만족한다고 하자. 그러면 $A\langle T\rangle$ 위에는
$\text{d}(T) = f$를 만족하는 유일한 미분 $\text{d}$가 존재하며,
이 $\text{d}$는 $A\langle T \rangle$의 분할 거듭제곱 구조와 양립한다.
\end{lemma}

\begin{proof}
직접 계산으로 증명할 수 있으며, 그 계산은 생략한다.
\end{proof}

\noindent
보조정리 \ref{dpa-lemma-compute-cohomology-adjoin-variable}에서는 몇 가지
특별한 경우에 $A\langle T \rangle$의 호몰로지를 계산할 것이다.
다음은 Avramov와 Halperin이 확장한 Tate의 구성이다.

\begin{lemma}
\label{dpa-lemma-tate-resolution}
$R \to S$를 가환환의 준동형이라 하자. 다음 분해가 존재한다.
$$
R \to A \to S
$$
이는 다음 성질들을 갖는다.
\begin{enumerate}
\item $(A, \text{d}, \gamma)$는 정의
\ref{dpa-definition-divided-powers-dga}에서와 같다.
\item $S$에 영 미분을 주면 $A \to S$는 유사동형이다.
\item $A_0 = R[x_j: j\in J] \to S$는 다항식 환에서 $S$로 가는
임의의 전사이다.
\item $A$는 $R$ 위의 등급 분할 거듭제곱 다항식 대수이다.
\end{enumerate}
마지막 조건은 Example \ref{dpa-example-adjoining-odd} 또는
\ref{dpa-example-adjoining-even}에서처럼 각 차수 $> 0$마다 변수 $T$의
집합을 차례로 첨가하여 $A_0$에서 $A$를 구성한다는 뜻이다.
더욱이 $R$이 뇌터이고 $R\to S$가 유한형이면, 각 차수에서
$A$의 생성원을 유한 개만 택할 수 있다.
\end{lemma}

\begin{proof}
$R$이 뇌터이고 $R\to S$가 유한형인 경우의 구성을 적어 보겠다.
이 가정들이 없어도 증명은 같지만, 각 차수에서 어떤 생성원 집합
(무한일 수도 있다)을 사용해야 한다.

\medskip\noindent
구성의 시작: $A(0) = R[x_1, \ldots, x_n]$을 (보통의) 다항식 환이라 하고
$A(0) \to S$를 전사라 하자. 등급은 $A(0)_0 = A(0)$ 및
$d \not = 0$일 때 $A(0)_d = 0$으로 준다.
따라서 $\text{d} = 0$이고, $n > 0$일 때 $\gamma_n$도 영이다.

\medskip\noindent
주어진 사상 $A(0) = R[x_1, \ldots, x_n] \to S$의 핵의 생성원
$f_1, \ldots, f_m \in R[x_1, \ldots, x_n]$을 택하자.
Example \ref{dpa-example-adjoining-odd}를 $m$번 적용하여
$$
A(1) = A(0)\langle T_1, \ldots, T_m\rangle
$$
을 얻는다. 이는 $\deg(T_i) = 1$인 등급 분할 거듭제곱 다항식 대수이다.
$\text{d}(T_i) = f_i$로 놓는다. $A(1)$은 $A(0)$ 위의
분할 거듭제곱 다항식 대수이고 $\text{d}(f_i) = 0$이므로,
보조정리 \ref{dpa-lemma-extend-differential}에 의해 이는 $A(1)$ 위의
미분으로 유일하게 연장된다.

\medskip\noindent
귀납 가정: 다음과 같은 분해들이 주어졌다고 하자.
$$
R \to A(0) \to A(1) \to \ldots \to A(m) \to S
$$
여기서 $A(0)$과 $A(1)$은 위와 같고, $2 \leq m' \leq m$인 각
$R \to A(m') \to S$는 보조정리 명제의 성질 (1), (4)를 만족하며,
(2)는 $m' > i \geq 0$일 때 $H_i(A(m')) \to H_i(S)$가
동형이라는 조건으로 바꾼다. 기초 단계는 $m = 1$이다.

\medskip\noindent
귀납 단계: 귀납 가정과 같은 $R \to A(m) \to S$가 있다고 하자.
군 $H_m(A(m))$을 생각하자. 이는 $H_0(A(m)) = S$ 위의 가군이다.
실제로 이는 $A(m)_0 = R[x_1, \ldots, x_n]$ 위의 유한형 가군인
$A(m)_m$의 부분몫이다. 따라서 이 가군의 생성원으로 가는 유한 개의 원소
$$
e_1, \ldots, e_t \in \Ker(\text{d} : A(m)_m \to A(m)_{m - 1})
$$
를 택할 수 있다. Example \ref{dpa-example-adjoining-odd} 또는
\ref{dpa-example-adjoining-even}를 $t$번 적용하여
$$
A(m + 1) = A(m)\langle T_1, \ldots, T_t\rangle
$$
을 얻는다. 이는 $\deg(T_i) = m + 1$인 등급 분할 거듭제곱 대수이다.
$\text{d}(T_i) = e_i$로 놓는다. $A(m+1)$은 $A(m)$ 위의
분할 거듭제곱 다항식 대수이고 $\text{d}(e_i) = 0$이므로,
이는 분할 거듭제곱 구조와 양립하는 $A(m + 1)$ 위의 미분으로
유일하게 연장된다. 차수 $m + 1$ 이상에만 항들을 첨가했으므로
$i < m$에 대해 $H_i(A(m + 1)) = H_i(A(m))$이다.
또한 구성에 의해 $H_m(A(m + 1)) = 0$임이 명백하다.

\medskip\noindent
증명을 마치기 위해, 다음 사상열이 존재함을 보였다는 데 유의하자.
$$
R \to A(0) \to A(1) \to \ldots \to A(m) \to A(m + 1) \to \ldots \to S
$$
이제 $A = \colim A(m)$로 놓으면 증명이 끝난다.
\end{proof}

\begin{lemma}
\label{dpa-lemma-tate-resoluton-pseudo-coherent-ring-map}
$R \to S$를 유사코히런트 환 사상이라 하자(More on Algebra, Definition
\ref{more-algebra-definition-pseudo-coherent-perfect}). 그러면
Lemma \ref{dpa-lemma-tate-resolution}가 성립하며, $S$의 분해 $A$는
각 차수마다 유한 개의 생성원만 갖도록 택할 수 있다.
\end{lemma}

\begin{proof}
이는 Lemma \ref{dpa-lemma-tate-resolution}와 완전히 같은 방식으로 증명된다.
유일한 추가점은 $A(m) \to S$가 주어졌을 때
$H_m = H_m(A(m))$가 유한 $R[x_1, \ldots, x_m]$-가군임을
보여야 한다는 것이다(그러면 다음 단계에서는 유한 개의 변수만 첨가하면 된다).
다음 복합체를 생각하자.
$$
\ldots \to A(m)_{m - 1} \to A(m)_m \to A(m)_{m - 1} \to
\ldots \to A(m)_0 \to S \to 0
$$
$S$는 유사코히런트 $R[x_1, \ldots, x_n]$-가군이고
$A(m)_i$는 유한 자유 $R[x_1, \ldots, x_n]$-가군이므로,
이는 유사코히런트 복합체이다. More on Algebra, Lemma
\ref{more-algebra-lemma-complex-pseudo-coherent-modules}를 보라.
이 복합체는 (호몰로지) 차수 $> m$에서 완전하므로, More on Algebra,
Lemma \ref{more-algebra-lemma-finite-cohomology}에 의해
$H_m$은 유한 $R$-가군이다.
\end{proof}

\begin{lemma}
\label{dpa-lemma-uniqueness-tate-resolution}
$R$을 가환환이라 하자. $(A, \text{d}, \gamma)$와
$(B, \text{d}, \gamma)$가 Definition
\ref{dpa-definition-divided-powers-dga}에서와 같다고 하자.
$\overline{\varphi} : H_0(A) \to H_0(B)$를 $R$-대수 사상이라 하자.
다음을 가정하자.
\begin{enumerate}
\item $A$는 $R$ 위의 등급 분할 거듭제곱 다항식 대수이다.
\item $k > 0$일 때 $H_k(B) = 0$이다.
\end{enumerate}
그러면 $\overline{\varphi}$를 올리는, 분할 거듭제곱과 양립하는
미분 등급 $R$-대수 사상 $\varphi : A \to B$가 존재한다.
\end{lemma}

\begin{proof}
이 가정은 차수 영의 어떤 다항식 생성원 집합, 양의 홀수 차수의
외대수 생성원, 양의 짝수 차수의 분할 거듭제곱 생성원을 $R$에
차례로 첨가하여 $A$를 얻는다는 뜻이다. 따라서 $A$에는 여과
$R \subset A(0) \subset A(1) \subset \ldots$가 있으며,
$A(m + 1)$은 $A(m)$에 차수 $m + 1$의 적절한 유형의 생성원
(이를 간단히 ‘분할 거듭제곱 생성원’이라 부른다)을 첨가하여 얻는다.
특히 $A(0) \to H_0(A)$는 $R$ 위의 (보통) 다항식 대수에서
$H_0(A)$로 가는 전사이다. 따라서 $\overline{\varphi}$를
$R$-대수 사상 $\varphi(0) : A(0) \to B_0$로 올릴 수 있다.

\medskip\noindent
% P03-STACKS-E961
$A(1) = A(0)\langle T_j:j\in J\rangle$라 쓰자. 여기서 $J$는
차수 $1$의 분할 거듭제곱 변수 $T_j$들의 어떤 집합이다.
$f_j \in B_0$를 $f_j = \varphi(0)(\text{d}(T_j))$로 놓자.
$\text{d}T_j$가 $H_0(A)$에서 영으로 가므로 $f_j$는
$H_0(B)$에서 영으로 간다. 따라서 $\text{d}(b_j) = f_j$인
$b_j \in B_1$을 찾을 수 있다. Lemma
\ref{dpa-lemma-divided-power-polynomial-algebra}의 분할 거듭제곱
다항식 대수의 보편 성질에 의해, $T_j$를 $f_j$로 보내는
$\varphi(0)$의 올림 $\varphi(1) : A(1) \to B$를 얻는다.

\medskip\noindent
어떤 $m \geq 1$에 대해 $\varphi(m)$을 구성했다면, 완전히 같은
방식으로 $\varphi(m + 1) : A(m + 1) \to B$를 구성할 수 있다.
세부 사항은 생략한다.
\end{proof}

\begin{lemma}
\label{dpa-lemma-divided-powers-on-tor}
$R$을 가환환이라 하자. $S$와 $T$를 가환 $R$-대수라 하자.
그러면
$$
\operatorname{Tor}_*^R(S, T).
$$
위에는 분할 거듭제곱을 갖춘 엄격히 등급 가환 $R$-대수의
표준 구조가 있다.
\end{lemma}

\begin{proof}
위와 같은 분해 $R \to A \to S$를 택하자. $A \to S$가 유사동형이고
$A_d$가 자유 $R$-가군이므로, 미분 등급 대수 $B = A \otimes_R T$가
Lemma에 표시된 Tor 군을 계산함을 알 수 있다. 전사
$R[y_j:j\in J] \to T$를 택하자. 그러면 $B$는 미분 등급 대수
$A[y_j:j\in J]$의 몫이며, 후자의 호몰로지는 차수 $0$에 놓인다
(이는 $S[y_j:j\in J]$와 같다). Lemma
\ref{dpa-lemma-base-change-div}에 의해 미분 등급 대수 $B$와
$A[y_j:j\in J]$는 미분과 양립하는 분할 거듭제곱 구조를 갖는다.
따라서 Lemma \ref{dpa-lemma-dpdga-good}에 의해 $H(B)$ 위에 원하는
분할 거듭제곱 구조를 얻는다.

\medskip\noindent
이렇게 구성한 분할 거듭제곱 대수 구조는 $A$의 선택과 무관하다.
실제로 $A'$이 두 번째 선택이면, Lemma
\ref{dpa-lemma-uniqueness-tate-resolution}에 의해 모든 구조와
$S$로 가는 증대 사상을 보존하는 사상 $A \to A'$이 존재한다.
그러면 유도된 사상
% P03-STACKS-E962
$B = A \otimes_R T \to A' \otimes_R T' = B'$도 모든 구조를
보존하며 유사동형이다. 따라서 유도된 Tor 대수의 동형은 곱과
분할 거듭제곱과 양립한다.
\end{proof}





\section{완전교차에의 응용}
\label{dpa-section-application-ci}

\noindent
$R$을 환이라 하자. $(A, \text{d}, \gamma)$가 Definition
\ref{dpa-definition-divided-powers-dga}에서와 같다고 하자.
차수 $2$의 {\it 미분}이란 다음 성질을 갖는 $R$-선형 사상
$\theta : A \to A$를 말한다.
\begin{enumerate}
\item $\theta(A_d) \subset A_{d - 2}$,
\item $\theta(xy) = \theta(x)y + x\theta(y)$,
\item $\theta$는 $\text{d}$와 가환한다,
% P03-STACKS-E1014
\item 모든 $x \in A_{2d}$와 모든 $d$에 대해
$\theta(\gamma_n(x)) = \theta(x) \gamma_{n - 1}(x)$이다.
\end{enumerate}
다음 Lemma에서 미분을 하나 구성한다.

\begin{lemma}
\label{dpa-lemma-get-derivation}
$R$을 환이라 하자. $(A, \text{d}, \gamma)$가 Definition
\ref{dpa-definition-divided-powers-dga}에서와 같다고 하자.
$R' \to R$을 그 핵의 제곱이 영이고 한 원소 $f$로 생성되는
환의 전사라 하자. $A$가 $R$ 위의 등급 분할 거듭제곱 다항식
대수이고 각 차수마다 변수가 유한 개이면, 미분
$\theta : A/IA \to A/IA$를 얻는다. 여기서 $I$는 $R$에서
$f$의 소멸자이다.
\end{lemma}

\begin{proof}
$A$가 분할 거듭제곱 다항식 대수이므로,
$A = A' \otimes_{R'} R$가 되는 $R'$ 위의 분할 거듭제곱
다항식 대수 $A'$을 찾을 수 있다. 또한
% P03-STACKS-E1015
$\text{d}$를 다음을 만족하는 $A'$ 위의 $R$-선형 연산자
$\text{d}$로 올릴 수 있다.
\begin{enumerate}
\item  동차원 $x, y \in A'$에 대해
$\text{d}(xy) = \text{d}(x)y + (-1)^{\deg(x)}x \text{d}(y)$이고,
\item  $x \in A'_{even, +}$에 대해
$\text{d}(\gamma_n(x)) = \text{d}(x) \gamma_{n - 1}(x)$이다.
\end{enumerate}
세부 사항은 생략한다(힌트: 한 번에 변수 하나씩 진행하라).
그러나 $A'$ 위에서 $\text{d}^2$이 영일 필요는 없다.
$\text{d}^2$이 $A'$을 $fA' \cong A/IA$로 보냄은 명백하다.
따라서 $\text{d}^2$은 $fA'$을 소멸시키며 사상
$A \to A/IA$를 통해 인수분해된다. $\text{d}^2$이 $R$-선형이므로
원하는 사상 $\theta : A/IA \to A/IA$를 얻는다.
미분의 성질들은 즉시 확인된다.
\end{proof}

\begin{lemma}
\label{dpa-lemma-compute-theta}
가정과 표기는 Lemma \ref{dpa-lemma-get-derivation}에서와 같다고 하자.
어떤 $n$, $m$, $f_j \in R[x_1, \ldots, x_n]$에 대해
$S = H_0(A)$가
$R[x_1, \ldots, x_n]/(f_1, \ldots, f_m)$
와 동형이라고 하자. 또한 관계식
$$
\sum r_j f_j = 0
$$
이 주어졌다고 하자. 여기서 $r_j \in R[x_1, \ldots, x_n]$이다.
$r_j, f_j$의 올림 $r'_j, f'_j \in R'[x_1, \ldots, x_n]$을 택하자.
어떤 $g \in R/I[x_1, \ldots, x_n]$에 대해
$\sum r'_j f'_j = gf$라 쓰자. $H_1(A) = 0$이고 각 $r_j$의
모든 계수가 $I$에 속하면, $S/IS$에서 $\theta(\xi) = g$를
만족하는 원소 $\xi \in H_2(A/IA)$가 존재한다.
\end{lemma}

\begin{proof}
$A(0) \subset A(1) \subset A(2) \subset \ldots$를 $A$의 여과라 하자.
여기서 $A(m)$은 $A(m - 1)$에 차수 $m$의 분할 거듭제곱 변수를
첨가하여 얻는다. 그러면 $A(0)$은 $R$ 위의 다항식 대수이고
$R$-대수 전사 $A(0) \to S$를 갖는다. 따라서 $S$로 가는 증대
사상들을 올리는 사상
$$
\varphi : R[x_1, \ldots, x_n] \to A(0)
$$
을 택할 수 있다. 다음으로, 어떤 차수 $1$의 분할 거듭제곱 변수
$T_i$들에 대해 $A(1) = A(0)\langle T_1, \ldots, T_t \rangle$이다.
$H_0(A) = S$이므로, $\text{d}(\xi_j) = \varphi(f_j)$인
$\xi_j \in \sum A(0)T_i$를 택할 수 있다. 그러면
$$
\text{d}\left(\sum \varphi(r_j) \xi_j\right) =
\sum  \varphi(r_j) \varphi(f_j) =  \sum \varphi(r_jf_j) = 0
$$
이다. $H_1(A) = 0$이므로,
$\text{d}(\xi) = \sum \varphi(r_j) \xi_j$인
$\xi \in A_2$를 택할 수 있다. $r_j$의 계수가 $I$에 속하면
$\varphi(r_j)$도 그러하다. 이 경우 $\text{d}(\xi)$는
$A_1/IA_1$에서 사라지므로, $\xi$는 $H_2(A/IA)$에서 한 류를 정의한다.

\medskip\noindent
Lemma \ref{dpa-lemma-get-derivation}의 증명에서 $\theta$를 구성할 때는
$A(i)$를 $A'(i)$로 차례로 올리고 미분 $\text{d}$도 올린다.
$\varphi$를 $\varphi' : R'[x_1, \ldots, x_n] \to A'(0)$로 올리자.
다음으로 $A'(1) = A'(0)\langle T_1, \ldots, T_t\rangle$이다.
또한 $\xi_j$를 $\xi'_j \in \sum A'(0)T_i$로 올릴 수 있다.
그러면 어떤 $a_j \in A'(0)$에 대해
$\text{d}(\xi'_j) = \varphi'(f'_j) + f a_j$이다.
$\xi$의 올림 $\xi' \in A'_2$을 생각하자. 그러면
$$
\text{d}(\xi') = \sum \varphi'(r'_j)\xi'_j + \sum fb_iT_i
$$
이다.
% P03-STACKS-E1016
여기서 $b_i \in A(0)$이다. $\text{d}$를 다시 적용하면
$$
\theta(\xi) = \sum \varphi'(r'_j)\varphi'(f'_j) +
\sum f \varphi'(r'_j) a_j + \sum fb_i \text{d}(T_i)
$$
을 얻는다. 첫 번째 항이 원하는 것을 준다. 두 번째 항은
$r_j$의 계수들이 $I$에 속하여 $f$에 의해 소멸되므로 영이다.
세 번째 항은 $\text{d}(T_i)$가 영으로 가므로 $H_0$에서 영으로 간다.
\end{proof}

\noindent
다음 Lemma의 증명 방법은 Gulliksen에 의한 것으로 보인다.

\begin{lemma}
\label{dpa-lemma-not-finite-pd}
$R' \to R$을 그 핵의 제곱이 영이고 한 원소 $f$로 생성되는
뇌터 환의 전사라 하자.
$S = R[x_1, \ldots, x_n]/(f_1, \ldots, f_m)$라 하자.
$\sum r_j f_j = 0$을 $R[x_1, \ldots, x_n]$의 관계식이라 하자.
다음을 가정하자.
\begin{enumerate}
\item 각 $r_j$의 계수는 $R$에서 $f$의 소멸자 $I$에 속한다,
\item 어떤 올림 $r'_j, f'_j \in R'[x_1, \ldots, x_n]$에 대해
$\sum r'_j f'_j = gf$이고, $g$는 $S/IS$에서 멱영이 아니다.
\end{enumerate}
그러면 $S$는 $R$ 위에서 유한 Tor 차원을 갖지 않는다
(즉, $S$는 완전 $R$-대수가 아니다).
\end{lemma}

\begin{proof}
Lemma \ref{dpa-lemma-tate-resolution}에서와 같은 Tate 분해
$R \to A \to S$를 택하자. $\xi \in H_2(A/IA)$와
$\theta : A/IA \to A/IA$를 Lemma \ref{dpa-lemma-get-derivation}과
\ref{dpa-lemma-compute-theta}에서 얻은 원소와 미분이라 하자.
다음이 성립함에 유의하자.
$$
\theta^n(\gamma_n(\xi)) = g^n
$$
이는 $H_0(A/IA) = S/IS$에서의 등식이다.
% P03-STACKS-E1017
따라서 $g$가 $S/IS$에서 멱영이 아니면, 모든 $n > 0$에 대해
$\xi^n$은 $H_{2n}(A/IA)$에서 영이 아니다.
$H_{2n}(A/IA) = \text{Tor}^R_{2n}(S, R/I)$이므로 결론을 얻는다.
\end{proof}

\noindent
다음 결과는 \cite{Rodicio}에서 찾을 수 있다.

\begin{lemma}
\label{dpa-lemma-injective}
$(A, \mathfrak m)$을 뇌터 국소환이라 하자.
$I \subset J \subset A$를 진아이디얼이라 하자.
$A/J$가 $A/I$ 위에서 유한 Tor 차원을 가지면
$I/\mathfrak m I \to J/\mathfrak m J$는 단사이다.
\end{lemma}

\begin{proof}
$f \in I$를 $I/\mathfrak m I$에서의 상은 영이 아니지만
$J/\mathfrak m J$에서는 영으로 가는 원소라 하자.
$\mathfrak mI \subset I' \subset I$이고 $I/I'$이 $f$의 상으로
생성되도록 아이디얼 $I'$을 택할 수 있다.
$R = A/I$와 $R' = A/I'$로 놓자.
어떤 $a_j \in A$에 대해 $J = (a_1, \ldots, a_m)$라 하자.
그러면 어떤 $b_j \in \mathfrak m$에 대해
$f = \sum b_j a_j$이다. $r_j, f_j \in R$와
$r'_j, f'_j \in R'$를 각각 $b_j, a_j$의 상이라 하자.
그러면 Lemma \ref{dpa-lemma-not-finite-pd}의 상황에 놓인다
(그 Lemma의 아이디얼 $I$는 $\mathfrak m_R$과 같다).
따라서 Lemma가 증명되었다.
\end{proof}

\begin{lemma}
\label{dpa-lemma-regular-sequence}
$(A, \mathfrak m)$을 뇌터 국소환이라 하자.
$I \subset J \subset A$를 진아이디얼이라 하자. 다음을 가정하자.
\begin{enumerate}
\item $A/J$는 $A/I$ 위에서 유한 Tor 차원을 갖는다,
\item $J$는 정칙열로 생성된다.
\end{enumerate}
그러면 $I$는 정칙열로 생성되고 $J/I$도 정칙열로 생성된다.
\end{lemma}

\begin{proof}
Lemma \ref{dpa-lemma-injective}에 의해
$I/\mathfrak m I \to J/\mathfrak m J$가 단사임을 알 수 있다.
따라서 $s \leq r$와 $J$의 최소 생성계 $f_1, \ldots, f_r$를
택하여, $f_1, \ldots, f_s$가 $I$에 속하고 $I$의 최소 생성계를
이루게 할 수 있다. More on Algebra, Lemmas
\ref{more-algebra-lemma-independence-of-generators}와
\ref{more-algebra-lemma-noetherian-finite-all-equivalent}에 의해
$J$의 모든 최소 생성계가 정칙열이므로 결론이 따른다.
\end{proof}

\begin{lemma}
\label{dpa-lemma-perfect-map-ci}
$R \to S$를 뇌터 국소환들의 국소 환 사상이라 하자.
$I \subset R$와 $J \subset S$를 $IS \subset J$인 아이디얼이라 하자.
$R \to S$가 평탄하고 $S/\mathfrak m_RS$가 정칙이면, 다음은 동치이다.
\begin{enumerate}
\item $J$는 정칙열로 생성되고 $S/J$는 $R/I$ 위의 가군으로서
유한 Tor 차원을 갖는다,
\item $J$는 정칙열로 생성되고
$\text{Tor}^{R/I}_p(S/J, R/\mathfrak m_R)$가 영이 아닌 $p$는
유한 개뿐이다,
\item $I$는 정칙열로 생성되고 $J/IS$는 $S/IS$에서
정칙열로 생성된다.
\end{enumerate}
\end{lemma}

\begin{proof}
(3)이 성립하면 $J$는 정칙열로 생성된다. 예를 들어
More on Algebra, Lemmas
\ref{more-algebra-lemma-join-koszul-regular-sequences}와
\ref{more-algebra-lemma-noetherian-finite-all-equivalent}를 보라.
또한 (3)이 성립하면 Koszul 복합체가 $S/IS$ 위의 $S/J$의
유한 자유 분해가 되므로
% P03-STACKS-E1018
$S/J = (S/I)/(J/I)$는 $S/IS$ 위에서 유한 사영 차원을 갖는다.
$R/I \to S/IS$가 평탄하므로, More on Algebra, Lemma
\ref{more-algebra-lemma-flat-push-tor-amplitude}에 의해
$S/J$는 $R/I$ 위에서 유한 Tor 차원을 갖는다.
따라서 (3)은 (1)을 함의한다.

\medskip\noindent
(1) $\Rightarrow$ (2)는 자명하다. (2)를 가정하자.
More on Algebra, Lemma
\ref{more-algebra-lemma-perfect-over-regular-local-ring}에 의해
$S/J$는 $S/IS$ 위에서 유한 Tor 차원을 갖는다.
따라서 Lemma \ref{dpa-lemma-regular-sequence}를 적용하면
$IS$와 $J/IS$가 정칙열로 생성됨을 알 수 있다.
$f_1, \ldots, f_r \in I$를 $I$의 최소 생성계라 하자.
$R \to S$가 평탄하므로, $f_1, \ldots, f_r$는 $S$에서
$IS$의 최소 생성계를 이룬다. 따라서
$f_1, \ldots, f_r \in R$은 그 상들이 $S$에서 정칙열을 이루는
원소열이다. More on Algebra, Lemmas
\ref{more-algebra-lemma-independence-of-generators}와
\ref{more-algebra-lemma-noetherian-finite-all-equivalent}를 보라.
그러므로 Algebra, Lemma \ref{algebra-lemma-flat-increases-depth}에 의해
$f_1, \ldots, f_r$는 $R$에서 정칙열이다.
\end{proof}






\section{국소 완전교차환}
\label{dpa-section-lci}

\noindent
$(A, \mathfrak m)$을 뇌터 완비 국소환이라 하자.
Cohen 구조 정리(Algebra, Theorem
\ref{algebra-theorem-cohen-structure-theorem} 참조)에 의해
$A$를 어떤 정칙 뇌터 완비 국소환 $R$의 상환으로 쓸 수 있다.
$R$가 정칙 국소환이고 핵이 정칙열로 생성되는 어떤 전사
$R \to A$가 존재하면 $A$를 {\it 완전교차}라고 하자.
다음 Lemma는 이 개념이 전사의 선택과 무관함을 보인다.

\begin{lemma}
\label{dpa-lemma-ci-well-defined}
$(A, \mathfrak m)$을 뇌터 완비 국소환이라 하자. 다음은 동치이다.
\begin{enumerate}
\item $R$가 정칙 국소환인 모든 국소환의 전사 $R \to A$에 대해
$R \to A$의 핵은 정칙열로 생성된다,
\item $R$가 정칙 국소환인 어떤 국소환의 전사 $R \to A$에 대해
$R \to A$의 핵은 정칙열로 생성된다.
\end{enumerate}
\end{lemma}

\begin{proof}
$k$를 $A$의 잉여체라 하자. $k$의 표수가 $p > 0$이면,
% P03-STACKS-E1019
$\Lambda$를 잉여체가 $k$인 Cohen 환(Algebra, Definition
\ref{algebra-definition-cohen-ring} 및 Algebra, Lemma
\ref{algebra-lemma-cohen-rings-exist})이라 하자.
$k$의 표수가 $0$이면 $\Lambda = k$로 놓는다.
임의의 $n$에 대해 $\Lambda[[x_1, \ldots, x_n]]$은
$\mathfrak m$-진 위상에서 각각 $\mathbf{Z}$ 또는 $\mathbf{Q}$ 위에
형식적으로 매끄럽다는 것을 상기하자. More on Algebra, Lemma
\ref{more-algebra-lemma-power-series-ring-over-Cohen-fs}를 보라.
Cohen 구조 정리(Algebra, Theorem
\ref{algebra-theorem-cohen-structure-theorem})에서와 같은 전사
$\Lambda[[x_1, \ldots, x_n]] \to A$를 하나 고정하자.

\medskip\noindent
$R \to A$를 정칙 국소환 $R$에서 나오는 전사라 하자.
$f_1, \ldots, f_r$를 $\Ker(R \to A)$의 최소 생성계라 하자.
뇌터 국소환의 한 아이디얼이 정칙열로 생성될 필요충분조건은
그 임의의 최소 생성계가 정칙열인 것이라는 사실을 이하에서
별도의 언급 없이 사용한다. Algebra, Lemmas
\ref{algebra-lemma-flat-increases-depth}와
\ref{algebra-lemma-completion-flat}에 의해,
$f_1, \ldots, f_r$가 $R$에서 정칙열일 필요충분조건은
$f_1, \ldots, f_r$가 완비화 $R^\wedge$에서 정칙열인 것이다.
또한 다음이 성립한다.
% P03-STACKS-E963
$$
R^\wedge/(f_1, \ldots, f_r)R^\wedge =
(R/(f_1, \ldots, f_n))^\wedge = A^\wedge = A
$$
왜냐하면 $A$가 $\mathfrak m_A$-진 완비이기 때문이다
(첫 번째 등식은 Algebra, Lemma
\ref{algebra-lemma-completion-tensor}에 의한다).
마지막으로 $R$가 정칙이므로 환 $R^\wedge$도 정칙이다
(More on Algebra, Lemma
\ref{more-algebra-lemma-completion-regular}).
따라서 $R$가 완비라고 가정해도 된다.

\medskip\noindent
$R$가 완비이면 주어진 사상
$\Lambda[[x_1, \ldots, x_n]] \to A$를 올리는 사상
$\Lambda[[x_1, \ldots, x_n]] \to R$을 택할 수 있다.
More on Algebra, Lemma
\ref{more-algebra-lemma-lift-continuous}를 보라.
$R \to A$의 핵의 생성원들로 가는 변수
$y_1, \ldots, y_m$을 더 첨가하면
$\Lambda[[x_1, \ldots, x_n, y_1, \ldots, y_m]] \to R$이
전사라고 가정할 수 있다(몇몇 세부 사항은 생략한다).
그러면 다음 가환 도식을 생각할 수 있다.
$$
\xymatrix{
\Lambda[[x_1, \ldots, x_n, y_1, \ldots, y_m]] \ar[r] \ar[d] & R \ar[d] \\
\Lambda[[x_1, \ldots, x_n]] \ar[r] & A
}
$$
Algebra, Lemma \ref{algebra-lemma-ci-well-defined}에 의해,
$R \to A$에 대한 조건은 고정해 둔 사상\hfil\break
$\Lambda[[x_1, \ldots, x_n]] \to A$에 대한 조건과 동치이다.
이로써 Lemma의 증명이 끝난다.
\end{proof}

\noindent
다음 두 Lemma는 위 정의가 자연스러운지를 확인한다.

\begin{lemma}
\label{dpa-lemma-quotient-regular-ring-by-regular-sequence}
$R$을 정칙환이라 하고 $\mathfrak p \subset R$를 소 아이디얼이라 하자.
$f_1, \ldots, f_r \in \mathfrak p$를 정칙열이라 하자.
그러면
$$
A = (R/(f_1, \ldots, f_r))_\mathfrak p =
R_\mathfrak p/(f_1, \ldots, f_r)R_\mathfrak p
$$
의 완비화는 위에서 정의한 의미의 완전교차이다.
\end{lemma}

\begin{proof}
$A$의 완비화는
$A^\wedge = R_\mathfrak p^\wedge/(f_1, \ldots, f_r)R_\mathfrak p^\wedge$
와 같다. 뇌터 환 $R_\mathfrak p$ 위의 유한 가군에 대한 완비화가
완전함자이기 때문이다(Algebra, Lemma
\ref{algebra-lemma-completion-tensor}).
Algebra, Lemmas \ref{algebra-lemma-completion-flat}와
\ref{algebra-lemma-flat-increases-depth}에 의해
열 $f_1, \ldots, f_r$의 $R_\mathfrak p$에서의 상은 정칙열이다.
또한 More on Algebra, Lemma
\ref{more-algebra-lemma-completion-regular}에 의해
$R_\mathfrak p^\wedge$은 정칙 국소환이다.
따라서 완비 국소환에 대한 완전교차의 정의로부터 결론이 따른다.
\end{proof}

\noindent
다음 Lemma는 Algebra, Lemma \ref{algebra-lemma-lci}의 유사이다.

\begin{lemma}
\label{dpa-lemma-quotient-regular-ring}
$R$을 정칙환이라 하고 $\mathfrak p \subset R$를 소 아이디얼이라 하자.
$I \subset \mathfrak p$를 아이디얼이라 하자.
$A = (R/I)_\mathfrak p = R_\mathfrak p/I_\mathfrak p$로 놓자.
다음은 동치이다.
\begin{enumerate}
\item $A$의 완비화는 위의 의미에서 완전교차이다,
\item $I_\mathfrak p \subset R_\mathfrak p$는 정칙열로 생성된다,
\item 가군 $(I/I^2)_\mathfrak p$는
$\dim(R_\mathfrak p) - \dim(A)$개의 원소로 생성될 수 있다,
% P03-STACKS-E597
\item 여기에 더 추가할 것.
\end{enumerate}
\end{lemma}

\begin{proof}
$R$을 $\mathfrak p$에서의 국소화로 바꾸어도 되며 실제로 그렇게 한다.
그러면 $\mathfrak p = \mathfrak m$은 $R$의 극대 아이디얼이고
$A = R/I$이다. $f_1, \ldots, f_r \in I$를 최소 생성계라 하자.
$A$의 완비화는
$A^\wedge = R^\wedge/(f_1, \ldots, f_r)R^\wedge$
와 같다. 뇌터 환 $R_\mathfrak p$ 위의 유한 가군에 대한 완비화가
완전함자이기 때문이다(Algebra, Lemma
\ref{algebra-lemma-completion-tensor}).

\medskip\noindent
(1)이 성립하면 가정에 의해 열 $f_1, \ldots, f_r$의 $R^\wedge$에서의
상은 정칙열이다. 따라서 Algebra, Lemmas
\ref{algebra-lemma-completion-flat}와
\ref{algebra-lemma-flat-increases-depth}에 의해 이는 $R$에서도
정칙열이다. 그러므로 (1)은 (2)를 함의한다.

\medskip\noindent
(3)이 성립한다고 하자. $c = \dim(R) - \dim(A)$로 놓고,
$I/I^2$의 생성원들로 가는 $f_1, \ldots, f_c \in I$를 택하자.
% P03-STACKS-E1020
Nakayama의 Lemma(Algebra, Lemma \ref{algebra-lemma-NAK})에 의해
$I = (f_1, \ldots, f_c)$임을 알 수 있다.
$R$는 정칙이고 따라서 Cohen--Macaulay이므로
(Algebra, Proposition \ref{algebra-proposition-CM-module}),
Algebra, Proposition \ref{algebra-proposition-CM-module}에 의해
$f_1, \ldots, f_c$는 정칙열이다. 따라서 (3)은 (2)를 함의한다.
마지막으로 Lemma
\ref{dpa-lemma-quotient-regular-ring-by-regular-sequence}에 의해
(2)는 (1)을 함의한다.
\end{proof}

\noindent
다음 결과는 Avramov에 의한다. \cite{Avramov}를 보라.

\begin{proposition}
\label{dpa-proposition-avramov}
$A \to B$를 뇌터 국소환들 사이의 평탄한 국소 준동형이라 하자.
다음은 동치이다.
\begin{enumerate}
\item $B^\wedge$은 완전교차이다,
\item $A^\wedge$과 $(B/\mathfrak m_A B)^\wedge$은 완전교차이다.
\end{enumerate}
\end{proposition}

\begin{proof}
다음 도식을 생각하자.
$$
\xymatrix{
B \ar[r] & B^\wedge \\
A \ar[u] \ar[r] & A^\wedge \ar[u]
}
$$
수평 사상들은 충실히 평탄하므로(Algebra, Lemma
\ref{algebra-lemma-completion-faithfully-flat}),
오른쪽 수직 화살표는 평탄하다(예를 들어 Algebra, Lemma
\ref{algebra-lemma-criterion-flatness-fibre-Noetherian}에 의한다).
또한
$(B/\mathfrak m_A B)^\wedge = B^\wedge/\mathfrak m_{A^\wedge} B^\wedge$
이 Algebra, Lemma \ref{algebra-lemma-completion-tensor}에 의해 성립한다.
따라서 $A$와 $B$가 완비 뇌터 국소환이라고 가정해도 된다.

\medskip\noindent
$A$와 $B$가 완비 뇌터 국소환이라고 하자.
More on Algebra, Lemma
\ref{more-algebra-lemma-embed-map-Noetherian-complete-local-rings}에서와
같은 도식
$$
\xymatrix{
S \ar[r] & B \\
R \ar[u] \ar[r] & A \ar[u]
}
$$
을 택하자. $I = \Ker(R \to A)$와 $J = \Ker(S \to B)$로 놓자.
$R/I = A \to B = S/J$가 평탄하므로 사상
% P03-STACKS-E1021
$J/IS \otimes_R R/\mathfrak m_R \to J/J \cap \mathfrak m_R S$
은 동형이다. 따라서 $J/IS$의 최소 생성계는
$\Ker(S/\mathfrak m_R S \to B/\mathfrak m_A B)$의 최소 생성계로 간다.
마지막으로 $S/\mathfrak m_R S$는 정칙 국소환이다.

\medskip\noindent
(1)이 성립한다고 하자. 즉, $J$는 정칙열로 생성된다고 하자.
$A = R/I \to B = S/J$가 평탄하므로
Lemma \ref{dpa-lemma-perfect-map-ci}를 적용할 수 있고,
$I$와 $J/IS$가 정칙열로 생성됨을 얻는다.
Algebra, Lemma
\ref{algebra-lemma-dimension-base-fibre-equals-total}에 의해
$\dim(B) = \dim(A) + \dim(B/\mathfrak m_A B)$ 및
$\dim(S/IS) = \dim(A) + \dim(S/\mathfrak m_R S)$
가 성립한다. 따라서 $J/IS$는
$$
\dim(S/IS) - \dim(S/J) = \dim(S/\mathfrak m_R S) - \dim(B/\mathfrak m_A B)
$$
개의 원소로 생성된다(Algebra, Lemma
\ref{algebra-lemma-one-equation}).
따라서 위에서 본 바와 같이
$\Ker(S/\mathfrak m_R S \to B/\mathfrak m_A B)$도
같은 수의 원소로 생성된다. 그러므로
$\Ker(S/\mathfrak m_R S \to B/\mathfrak m_A B)$은
정칙열로 생성된다.
예를 들어 Lemma \ref{dpa-lemma-quotient-regular-ring}를 보라.
이로써 (2)가 성립함을 알 수 있다.

\medskip\noindent
(2)가 성립하면 $I$와 $J/J \cap \mathfrak m_RS$는
% P03-STACKS-E1021
정칙열로 생성된다. 이 생성원들을 올리고(위 참조),
$R/I \to S/IS$의 평탄성과 Grothendieck의 Lemma
(Algebra, Lemma \ref{algebra-lemma-grothendieck-regular-sequence})를
사용하면 $J/IS$가 $S/IS$에서 정칙열로 생성됨을 알 수 있다.
따라서 Lemma \ref{dpa-lemma-perfect-map-ci}에 의해 $J$는 정칙열로
생성되고, 이에 따라 (1)이 성립한다.
\end{proof}

\begin{definition}
\label{dpa-definition-lci}
$A$를 뇌터 환이라 하자.
\begin{enumerate}
\item $A$가 국소환일 때 그 완비화가 위의 의미에서 완전교차이면
$A$를 {\it 완전교차}라고 한다.
\item 일반적으로 모든 국소환이 완전교차이면
$A$를 {\it 국소 완전교차}라고 한다.
\end{enumerate}
\end{definition}

\noindent
이것이 Algebra, Definitions \ref{algebra-definition-lci-field}와
\ref{algebra-definition-lci-local-ring}에서 도입한 용어와
충돌하지 않음을 아래에서 확인할 것이다.
그러나 먼저 $A$가 뇌터 국소 완전교차이면 $A$는 국소 완전교차,
즉 그 모든 국소환이 완전교차임을 보임으로써 이 정의가
``말이 됨''을 보인다.

\begin{lemma}
\label{dpa-lemma-ci-good}
$(A, \mathfrak m)$을 뇌터 국소환이라 하고
$\mathfrak p \subset A$를 소 아이디얼이라 하자.
$A$가 완전교차이면 $A_\mathfrak p$도 완전교차이다.
\end{lemma}

\begin{proof}
$\mathfrak p$ 위에 놓이는 $A^\wedge$의 소 아이디얼
$\mathfrak q$를 택하자(Algebra, Lemma
\ref{algebra-lemma-completion-faithfully-flat}에 의해
$A \to A^\wedge$이 충실히 평탄하므로 이는 가능하다).
그러면 $A_\mathfrak p \to (A^\wedge)_\mathfrak q$는
평탄한 국소 환 준동형이다. 따라서 Proposition
\ref{dpa-proposition-avramov}에 의해 $A_\mathfrak p$가 완전교차일
필요충분조건은 $(A^\wedge)_\mathfrak q$가 완전교차인 것이다.
그러므로 $A$가 완비인 경우에 Lemma를 증명하면 충분하다
(이것이 증명의 핵심 단계이다).

\medskip\noindent
$A$가 완비라고 하자. 정의에 의해 어떤 정칙 국소환 $R$의
정칙열 $f_1, \ldots, f_r$에 대해
$A = R/(f_1, \ldots, f_r)$로 쓸 수 있다.
$\mathfrak q \subset R$를 $\mathfrak p$에 대응하는 소 아이디얼이라
하자. $f_1, \ldots, f_r \in \mathfrak q$이고
$A_\mathfrak p = R_\mathfrak q/(f_1, \ldots, f_r)R_\mathfrak q$.
따라서 Lemma \ref{dpa-lemma-quotient-regular-ring-by-regular-sequence}에
의해 $A_\mathfrak p$는 완전교차이다.
\end{proof}

\begin{lemma}
\label{dpa-lemma-check-lci-at-maximal-ideals}
$A$를 뇌터 환이라 하자. $A$가 국소 완전교차일 필요충분조건은
$A$의 모든 극대 아이디얼 $\mathfrak m$에 대해
$A_\mathfrak m$이 완전교차인 것이다.
\end{lemma}

\begin{proof}
이는 Lemma \ref{dpa-lemma-ci-good}와 정의들로부터 바로 따른다.
\end{proof}

\begin{lemma}
\label{dpa-lemma-check-lci-agrees}
$S$를 체 $k$ 위의 유한형 대수라 하자.
\begin{enumerate}
\item 소 아이디얼 $\mathfrak q \subset S$에 대해 국소환
$S_\mathfrak q$가 Algebra, Definition
\ref{algebra-definition-lci-local-ring}의 의미에서 완전교차일
필요충분조건은 $S_\mathfrak q$가 Definition
\ref{dpa-definition-lci}의 의미에서 완전교차인 것이다,
\item $S$가 Algebra, Definition
\ref{algebra-definition-lci-field}의 의미에서 국소 완전교차일
필요충분조건은 $S$가 Definition \ref{dpa-definition-lci}의 의미에서
국소 완전교차인 것이다.
\end{enumerate}
\end{lemma}

\begin{proof}
(1)을 증명하자. $k[x_1, \ldots, x_n] \to S$를 전사라 하자.
$\mathfrak p \subset k[x_1, \ldots, x_n]$를
$\mathfrak q$에 대응하는 소 아이디얼이라 하자.
$I \subset k[x_1, \ldots, x_n]$를 이 전사의 핵이라 하자.
$k[x_1, \ldots, x_n]_\mathfrak p \to S_\mathfrak q$는
핵이 $I_\mathfrak p$인 전사이다.
Algebra, Proposition
\ref{algebra-proposition-finite-gl-dim-polynomial-ring}에 의해
$k[x_1, \ldots, x_n]$은 정칙환이다. 따라서 Lemma
\ref{dpa-lemma-quotient-regular-ring}과 Algebra, Lemma
\ref{algebra-lemma-lci-local}을 결합하면 (1)의 두 개념이
동치임을 얻는다.

\medskip\noindent
(1)을 증명했으므로 (2)의 동치는 정의와 Algebra, Lemma
\ref{algebra-lemma-lci-global}에서 따른다.
\end{proof}

\begin{lemma}
\label{dpa-lemma-avramov}
$A \to B$를 뇌터 국소환들 사이의 평탄한 국소 준동형이라 하자.
다음은 동치이다.
\begin{enumerate}
\item $B$는 완전교차이다,
\item $A$와 $B/\mathfrak m_A B$는 완전교차이다.
\end{enumerate}
\end{lemma}

\begin{proof}
이제 정의가 말이 됨을 알았으므로, 이는 비자명한 Proposition
\ref{dpa-proposition-avramov}의 자명한 재정식화이다.
\end{proof}









\section{국소 완전교차 준동형}
\label{dpa-section-lci-homomorphisms}

\noindent
$A \to B$를 뇌터 완비 국소환들 사이의 국소 준동형이라 하자.
Cohen 구조 정리의 한 귀결로서 다음과 같은 뇌터 완비 국소환들의
가환 도식을 찾을 수 있다.
$$
\xymatrix{
S \ar[r] & B \\
& A \ar[lu] \ar[u]
}
$$
여기서 $S \to B$는 전사이고, $A \to S$는 평탄하며,
$S/\mathfrak m_A S$는 정칙 국소환이다. 이는 More on Algebra, Lemma
\ref{more-algebra-lemma-embed-map-Noetherian-complete-local-rings}에서
따른다. $S$가 뇌터 국소환이고, $A \to S \to B$가 국소 환
준동형들이며, $S \to B$가 전사이고, $A \to S$가 평탄하며,
$S/\mathfrak m_AS$가 정칙이면 (잠정적으로)
$A \to S \to B$를 $A \to B$의 {\it 좋은 인수분해}라고 하자.
$S \to B$의 핵이 정칙열로 생성되는 어떤 좋은 인수분해
$A \to S \to B$가 존재하면 $A \to B$를
{\it 완전교차 준동형}이라고 하자.
다음 Lemma는 이 개념이 도식의 선택과 무관함을 보인다.

\begin{lemma}
\label{dpa-lemma-ci-map-well-defined}
$A \to B$를 뇌터 완비 국소환들 사이의 국소 준동형이라 하자.
다음은 동치이다.
\begin{enumerate}
\item 어떤 좋은 인수분해 $A \to S \to B$에 대해
$S \to B$의 핵은 정칙열로 생성된다,
\item 모든 좋은 인수분해 $A \to S \to B$에 대해
$S \to B$의 핵은 정칙열로 생성된다.
\end{enumerate}
\end{lemma}

\begin{proof}
$A \to S \to B$를 좋은 인수분해라 하자.
$B$가 완비이므로 $S^\wedge$을 $S$의 완비화라 할 때
인수분해 $A \to S^\wedge \to B$를 얻는다.
이것도 좋은 인수분해임에 유의하자. 환 준동형
$S \to S^\wedge$은 평탄하므로(Algebra, Lemma
\ref{algebra-lemma-completion-flat}), $A \to S^\wedge$은 평탄하다.
$S/\mathfrak m_A S$가 정칙이므로 환
$S^\wedge/\mathfrak m_A S^\wedge = (S/\mathfrak m_A S)^\wedge$도
정칙이다(More on Algebra, Lemma
\ref{more-algebra-lemma-completion-regular}).
$f_1, \ldots, f_r$를 $\Ker(S \to B)$의 최소 생성계라 하자.
뇌터 국소환의 한 아이디얼이 정칙열로 생성될 필요충분조건은
그 임의의 최소 생성계가 정칙열인 것이라는 사실을 이하에서
별도의 언급 없이 사용한다. Algebra, Lemma
\ref{algebra-lemma-flat-increases-depth}에 의해
$f_1, \ldots, f_r$가 $S$에서 정칙열일 필요충분조건은
$f_1, \ldots, f_r$가 완비화 $S^\wedge$에서 정칙열인 것이다.
또한 다음이 성립한다.
% P03-STACKS-E964
% P03-STACKS-E965
$$
S^\wedge/(f_1, \ldots, f_r)R^\wedge =
(S/(f_1, \ldots, f_n))^\wedge = B^\wedge = B
$$
$B$가 $\mathfrak m_B$-진 완비이기 때문이다(첫 번째 등식은
Algebra, Lemma \ref{algebra-lemma-completion-tensor}에 의한다).
따라서 $S \to B$의 핵이 정칙열로 생성될 필요충분조건은
$S^\wedge \to B$의 핵이 정칙열로 생성되는 것이다.
그러므로 $S$가 완비인 좋은 인수분해들만 생각하면 충분하다.

\medskip\noindent
두 인수분해 $A \to S \to B$와 $A \to S' \to B$가 있고
$S$와 $S'$가 완비라고 하자. More on Algebra, Lemma
\ref{more-algebra-lemma-dominate-two-surjections}에 의해
환 $S \times_B S'$는 뇌터 완비 국소환이다.
따라서 More on Algebra, Lemma
\ref{more-algebra-lemma-embed-map-Noetherian-complete-local-rings}를
사용하여 $S''$가 완비인 좋은 인수분해
$A \to S'' \to S \times_B S'$를 택할 수 있다.
그러므로 다음을 보이면 충분하다. $A \to S' \to S \to B$가
비교 가능한 좋은 인수분해들이면, $\Ker(S \to B)$가 정칙열로
생성될 필요충분조건은 $\Ker(S' \to B)$가 정칙열로 생성되는 것이다.

\medskip\noindent
% P03-STACKS-E966
$A \to S' \to S \to B$를 비교 가능한 좋은 인수분해들이라 하자.
먼저 $S'/\mathfrak m_R S' \to S/\mathfrak m_R S$가 정칙 국소환들의
전사이므로 그 핵은 정칙열
$\overline{x}_1, \ldots, \overline{x}_c \in
\mathfrak m_{S'}/\mathfrak m_R S'$
로 생성되며, 이는 정칙 국소환 $S'/\mathfrak m_R S'$의
정칙 매개변수계로 연장될 수 있다(Algebra, Lemma
\ref{algebra-lemma-regular-quotient-regular} 참조).
$I = \Ker(S' \to S)$로 놓자. $S$의 $R$ 위 평탄성에 의해
$$
I/\mathfrak m_R I =
\Ker(S'/\mathfrak m_R S' \to S/\mathfrak m_R S) =
(\overline{x}_1, \ldots, \overline{x}_c).
$$
올림 $x_1, \ldots, x_c \in I$를 택하자.
$S'$가 $R$ 위에서 평탄하므로 이 올림들은 $I$를 생성하는
정칙열을 이룬다. Algebra, Lemma
\ref{algebra-lemma-grothendieck-regular-sequence}를 보라.

\medskip\noindent
따라서 $\Ker(S \to B)$도 정칙열로 생성되면
$\Ker(S' \to B)$도 그러하다. More on Algebra, Lemmas
\ref{more-algebra-lemma-join-koszul-regular-sequences}와
\ref{more-algebra-lemma-noetherian-finite-all-equivalent}를 보라.

\medskip\noindent
역으로 $J = \Ker(S' \to B)$가 정칙열로 생성된다고 하자.
$I$의 생성원 $x_1, \ldots, x_c$가
$\mathfrak m_{S'}/\mathfrak m_{S'}^2$의 선형 독립인 원소들로 가므로
$I/\mathfrak m_{S'}I \to J/\mathfrak m_{S'}J$는 단사이다.
따라서 $J$의 최소 생성계
$x_1, \ldots, x_c, y_1, \ldots, y_d$가 존재한다.
그러면 $x_1, \ldots, x_c, y_1, \ldots, y_d$는 정칙열이고,
$y_1, \ldots, y_d$의 $S$에서의 상들은
$\Ker(S \to B)$를 생성하는 정칙열을 이룬다.
이로써 Lemma의 증명이 끝난다.
\end{proof}

\noindent
% P03-STACKS-E1022
다음 Proposition에서 Tor의 소멸에 관한 조건은 특히 $B$가 $A$ 위에서
유한 Tor 차원을 가지는 경우, 따라서 특히 $B$가 $A$ 위에서 평탄한
경우에 성립함에 유의하자.

\begin{proposition}
\label{dpa-proposition-avramov-map}
$A \to B$를 뇌터 국소환들 사이의 국소 준동형이라 하자.
다음은 동치이다.
\begin{enumerate}
\item $B$는 완전교차이고
$\text{Tor}^A_p(B, A/\mathfrak m_A)$가 영이 아닌 $p$는 유한 개뿐이다,
\item $A$는 완전교차이고 $A^\wedge \to B^\wedge$은
위에서 정의한 의미의 완전교차 준동형이다.
\end{enumerate}
\end{proposition}

\begin{proof}
$F_\bullet \to A/\mathfrak m_A$를 유한 자유 $A$-가군들에 의한
분해라 하자.
$\text{Tor}^A_p(B, A/\mathfrak m_A)$
가 복합체 $F_\bullet \otimes_A B$의 제 $p$ 호몰로지임에 유의하자.
$F_\bullet^\wedge = F_\bullet \otimes_A A^\wedge$을 완비화라 하자.
그러면 $F_\bullet^\wedge$은 유한 자유 $A^\wedge$-가군들에 의한
$A^\wedge/\mathfrak m_{A^\wedge}$의 분해이다
($A \to A^\wedge$이 평탄하고 유한 가군에 대한 완비화가
완전함자이기 때문이다. Algebra, Lemmas
\ref{algebra-lemma-completion-tensor}와
\ref{algebra-lemma-completion-flat}을 보라).
따라서
$$
F_\bullet^\wedge \otimes_{A^\wedge} B^\wedge =
F_\bullet \otimes_A B \otimes_B B^\wedge
$$
이고, $B \to B^\wedge$의 평탄성에 의해
$$
\text{Tor}^{A^\wedge}_p(B^\wedge, A^\wedge/\mathfrak m_{A^\wedge}) =
\text{Tor}^A_p(B, A/\mathfrak m_A) \otimes_B B^\wedge
$$
를 얻는다. 이로써 국소 환 사상 $A \to B$에 관한 (1)의 조건은
국소 환 사상 $A^\wedge \to B^\wedge$에 관한 같은 조건과
동치임을 알 수 있다. 따라서 $A$와 $B$가 완비 뇌터 국소환이라고
가정해도 된다(나머지 조건들은 어느 경우든 완비화를 사용하여
정식화되어 있기 때문이다).

\medskip\noindent
$A$와 $B$가 완비 뇌터 국소환이라고 하자.
More on Algebra, Lemma
\ref{more-algebra-lemma-embed-map-Noetherian-complete-local-rings}에서와
같은 도식
$$
\xymatrix{
S \ar[r] & B \\
R \ar[u] \ar[r] & A \ar[u]
}
$$
을 택하자. $I = \Ker(R \to A)$와 $J = \Ker(S \to B)$로 놓자.
이제 Proposition은 Lemma \ref{dpa-lemma-perfect-map-ci}에서 따른다.
\end{proof}

\begin{remark}
\label{dpa-remark-no-good-ci-map}
% P03-STACKS-E967
일반 뇌터 환들 사이의 사상에 대해 ``국소 완전교차 준동형''의
좋은 개념을 정의하기는 어려워 보인다. 그 이유는 뇌터 국소환
$A$에 대해 $A \to A^\wedge$의 올들이 국소 완전교차환이 아니기
때문이다. 따라서 $A \to B$가 뇌터 국소환들 사이의 국소
준동형이고 완비화들의 사상 $A^\wedge \to B^\wedge$이 위에서
정의한 의미의 완전교차 준동형이더라도, 일반적으로
% P03-STACKS-E1023
$(A_\mathfrak p)^\wedge \to (B_\mathfrak q)^\wedge$은 위에서
정의한 의미의 완전교차 준동형이 {\bf 아니다}.
우수 뇌터 환들만을 다루면 해법을 얻을 수 있다.
더 일반적으로는 형식적 올들이 완전교차인 뇌터 환들을
다룰 수 있다. \cite{Rodicio-ci}를 보라.
이 이론은 Dualizing Complexes, Section
\ref{dualizing-section-formal-fibres}에서 전개한다.
\end{remark}

\noindent
% P03-STACKS-E968
이 절을 마치기 위해 위에서 정의한 개념을
More on Algebra, Section \ref{dpa-section-lci}에서 도입한 개념과 비교한다.

\begin{lemma}
\label{dpa-lemma-well-defined-if-you-can-find-good-factorization}
다음 뇌터 국소환들의 가환 도식을 생각하자.
$$
\xymatrix{
S \ar[r] & B \\
& A \ar[lu] \ar[u]
}
$$
여기서 $S \to B$는 전사이고, $A \to S$는 평탄하며,
$S/\mathfrak m_A S$는 정칙 국소환이다. 다음은 동치이다.
\begin{enumerate}
\item $\Ker(S \to B)$는 정칙열로 생성된다,
\item $A^\wedge \to B^\wedge$은 위에서 정의한 완전교차 준동형이다.
\end{enumerate}
\end{lemma}

\begin{proof}
생략한다. 힌트: 증명은 Lemma \ref{dpa-lemma-ci-map-well-defined}의
증명 첫 문단에서 한 논증과 동일하다.
\end{proof}

\begin{lemma}
\label{dpa-lemma-finite-type-lci-map}
$A$를 뇌터 환이라 하자. $A \to B$를 유한형 환 준동형이라 하자.
다음은 동치이다.
\begin{enumerate}
\item $A \to B$는 More on Algebra, Definition
\ref{more-algebra-definition-local-complete-intersection}의 의미에서
국소 완전교차이다,
% P03-STACKS-E1024
\item 모든 소 아이디얼 $\mathfrak q \subset B$에 대해
$\mathfrak p = A \cap \mathfrak q$로 놓으면, 환 준동형
$(A_\mathfrak p)^\wedge \to (B_\mathfrak q)^\wedge$은
위에서 정의한 의미의 완전교차 준동형이다.
\end{enumerate}
\end{lemma}

\begin{proof}
전사 $R = A[x_1, \ldots, x_n] \to B$를 택하자.
$A \to R$가 평탄하고 올들이 정칙임에 유의하자.
$I$를 $R \to B$의 핵이라 하자. (2)를 가정하자.
Lemma \ref{dpa-lemma-well-defined-if-you-can-find-good-factorization}와
Algebra, Lemma
\ref{algebra-lemma-regular-sequence-in-neighbourhood}에 의해
$I$는 국소적으로 정칙열로 생성됨을 알 수 있다.
달리 말해 (1)이 성립한다. 역으로 (1)을 가정하자.
$R$와 $B$를 국소화한 뒤에는 $I$가 Koszul 정칙열로 생성된다고
가정할 수 있다. More on Algebra, Lemma
\ref{more-algebra-lemma-noetherian-finite-all-equivalent}에 의해
$I$가 국소적으로 정칙열로 생성됨을 알 수 있다.
% P03-STACKS-E1025
따라서 Lemma
\ref{dpa-lemma-well-defined-if-you-can-find-good-factorization}에 의해
(2)가 성립한다. 몇몇 세부 사항은 생략한다.
\end{proof}

\begin{lemma}
\label{dpa-lemma-avramov-map-finite-type}
$A$를 뇌터 환이라 하자. $A \to B$를 그 상
$\Spec(B) \to \Spec(A)$이 $\Spec(A)$의 모든 닫힌점을 포함하는
유한형 환 준동형이라 하자. 다음은 동치이다.
\begin{enumerate}
\item $B$는 완전교차이고 $A \to B$는 유한 Tor 차원을 갖는다,
\item $A$는 완전교차이고 $A \to B$는 More on Algebra, Definition
\ref{more-algebra-definition-local-complete-intersection}의 의미에서
국소 완전교차이다.
\end{enumerate}
\end{lemma}

\begin{proof}
이는 Proposition \ref{dpa-proposition-avramov-map}을 Lemma
\ref{dpa-lemma-finite-type-lci-map}을 통해 재정식화한 것이다.
세부 사항은 생략한다.
\end{proof}








\section{매끄러운 환 준동형과 대각 사상}
\label{dpa-section-smooth-diagonal-perfect}

\noindent
이 절에서는 위의 내용을 사용하여 매끄러운 환 준동형을 그 대각 사상이
완전한 것들로 특징짓는다.

\begin{lemma}
\label{dpa-lemma-local-perfect-diagonal}
$A \to B$를 뇌터 국소환들 사이의 국소 환 준동형이라 하자.
$B$가 $A$ 위에서 평탄하고 본질적으로 유한형이라고 하자. 만일
$$
B \otimes_A B \longrightarrow B
$$
가 완전 환 준동형, 즉 $B$가 $B \otimes_A B$ 위에서 유한 Tor 차원을
가지면, $B$는 어떤 매끄러운 $A$-대수의 국소화이다.
\end{lemma}

\begin{proof}
$B$가 $A$ 위에서 본질적으로 유한형이므로 $B \otimes_A B$도 그러하며,
특히 $B \otimes_A B$는 뇌터 환이다. 따라서 $B \otimes_A B$의 상환
$B$는 $B \otimes_A B$ 위에서 유사코히런트이다(More on Algebra,
Lemma \ref{more-algebra-lemma-Noetherian-pseudo-coherent}). 이는 환
준동형의 완전성(More on Algebra, Definition
\ref{more-algebra-definition-pseudo-coherent-perfect})이 유한 Tor 차원
조건과 일치하는 이유를 설명한다.

\medskip\noindent
$B = R/K$로 쓸 수 있다. 여기서 $R$은 $A[x_1, \ldots, x_n]$을 어떤
소 아이디얼에서 국소화한 것이고 $K \subset R$은 아이디얼이다.
% P03-STACKS-E1026
$\mathfrak m \subset R \otimes_A R$을 전사
$R \otimes_A R \to B \otimes_A B \to B$에 의한 $B$의 극대 아이디얼의
역상인 극대 아이디얼이라 하자. 그러면 다음 전사들이 있다.
$$
(R \otimes_A R)_\mathfrak m \to (B \otimes_A B)_\mathfrak m \to B
$$
따라서 Lemma \ref{dpa-lemma-injective}에서와 같은 아이디얼
$I \subset J \subset (R \otimes_A R)_\mathfrak m$을 얻는다. 이로부터
$I/\mathfrak m I \to J/\mathfrak m J$가 단사임을 알 수 있다.

\medskip\noindent
$K = (f_1, \ldots, f_r)$이고 $r$이 최소라고 하자. $f_i \in R$은
$A[x_1, \ldots, x_n]$의 어떤 원소의 상이라고 가정해도 되고 실제로
그렇게 하며, 그 원소도 $f_i$로 나타낸다. $I$는
$f_1 \otimes 1, \ldots, f_r \otimes 1$과
$1 \otimes f_1, \ldots, 1 \otimes f_r$로 생성됨에 유의하자.
우리는 이것이 $I$의 최소 생성계라고 주장한다. 실제로 $\kappa$가
$R$, $B$, $(R \otimes_A R)_\mathfrak m$, $(B \otimes_A B)_\mathfrak m$의
공통 잉여체이면 다음 사상이 있다.
$R \otimes_A R \to R \otimes_A \kappa \oplus \kappa \otimes_A R$
이 사상은 $(R \otimes_A R)_\mathfrak m$을 통해 인수분해된다.
$B$가 $A$ 위에서 평탄하고 짧은 완전열
$0 \to K \to R \to B \to 0$이 있으므로
$K \otimes_A \kappa \subset R \otimes_A \kappa$임을 알 수 있다.
Algebra, Lemma \ref{algebra-lemma-flat-tor-zero}를 보라. 따라서 사상
$(R \otimes_A R)_\mathfrak m \to R \otimes_A \kappa \oplus \kappa \otimes_A R$
을 $I$에 제한하면 다음 사상을 얻는다.
% P03-STACKS-E1027
$$
I \to K \otimes_A \kappa \oplus \kappa \otimes_A K \to
K \otimes_B \kappa \oplus \kappa \otimes_B K.
$$
원소들
$f_1 \otimes 1, \ldots, f_r \otimes 1, 1 \otimes f_1, \ldots, 1 \otimes f_r$
은 이 사상의 공역의 한 기저로 간다. 이는 Nakayama의 Lemma
(Algebra, Lemma \ref{algebra-lemma-NAK})에 의해
$f_1, \ldots, f_r$가 $K \otimes_B \kappa$의 한 기저로 가기 때문이다.
이로써 주장이 증명되었다.

\medskip\noindent
아이디얼 $J$는 $f_1 \otimes 1, \ldots, f_r \otimes 1$과 원소들
$x_1 \otimes 1 - 1 \otimes x_1, \ldots,
x_n \otimes 1 - 1 \otimes x_n$로 생성된다(증명에는 이 원소들이
아이디얼 $J$에 포함됨을 아는 것으로 충분하다). 이제 다음과 같이 쓸 수 있다.
$$
f_i \otimes 1 - 1 \otimes f_i =
\sum g_{ij} (x_j \otimes 1 - 1 \otimes x_j)
$$
여기서 $g_{ij}$는 $(R \otimes_A R)_\mathfrak m$의 어떤 원소이다.
이는 $A[x_1, \ldots, x_n]$의 원소에 관한 일반적 사실이며 증명은 생략한다.
$a_{ij} \in \kappa$를 $g_{ij}$의 상이라 하자. 다른 계산에 의해
$a_{ij}$는 $\partial f_i / \partial x_j$의 $\kappa$에서의 상이다.
$I/\mathfrak m I \to J/\mathfrak m J$의 단사성과 위의 관찰들로부터
행렬 $(a_{ij})$가 최대 계수 $r$를 가져야 함을 알 수 있다. 다음과 같이 놓자.
$$
C = A[x_1, \ldots, x_n]/(f_1, \ldots, f_r)
$$
그리고 소박한 여접 복합체
$\NL_{C/A} \cong (C^{\oplus r} \to C^{\oplus n})$를 생각하자.
여기서 사상은 편미분들의 행렬로 주어진다.
% P03-STACKS-E1028
따라서 $\NL_{C/A} \otimes_A B$는 차수 $0$에 놓인 계수 $n - r$의
자유 $B$-가군과 동형이다. 그러므로 $B$에서 단원으로 가는 어떤
$g \in C$에 대해 $C_g$는 $A$ 위에서 매끄럽다. Algebra, Lemma
\ref{algebra-lemma-smooth-at-point}를 보라. 이로써 증명이 끝난다.
\end{proof}

\begin{lemma}
\label{dpa-lemma-perfect-diagonal}
$A \to B$를 뇌터 환들의 평탄한 유한형 환 준동형이라 하자. 만일
$$
B \otimes_A B \longrightarrow B
$$
가 완전 환 준동형, 즉 $B$가 $B \otimes_A B$ 위에서 유한 Tor 차원을
가지면, $B$는 매끄러운 $A$-대수이다.
\end{lemma}

\begin{proof}
이는 Lemma \ref{dpa-lemma-local-perfect-diagonal}와 매끄러운 환 준동형에
관한 일반적 사실들에서 따른다. Algebra, Lemmas
\ref{algebra-lemma-smooth-at-point}와
\ref{algebra-lemma-locally-smooth}를 보라. 또는 독자는 Lemma
\ref{dpa-lemma-local-perfect-diagonal}의 증명을 약간 수정하여 이 Lemma를
증명할 수 있다.
\end{proof}






\section{코노멀 가군의 자유성}
\label{dpa-section-freeness-conormal}

\noindent
Tate 분해와 그 위의 미분을 사용하면 이 절의 결과들의 (더 강한)
형태를 증명할 수 있다. \cite{Iyengar}를 보라. 더 초등적인 두 참고문헌은
\cite{Vasconcelos}와 \cite{Ferrand-lci}이다.

\begin{lemma}
\label{dpa-lemma-free-summand-in-ideal-finite-proj-dim}
\begin{reference}
\cite{Vasconcelos}
\end{reference}
$R$을 뇌터 국소환이라 하자. $I \subset R$을 $R$ 위에서 유한 사영
차원을 갖는 아이디얼이라 하자. $F \subset I/I^2$가 $R/I$와 동형인
직합인자이면, $x$의 $I/I^2$에서의 상이 $F$를 생성하는 비영인자
$x \in I$가 존재한다.
\end{lemma}

\begin{proof}
가정에 의해 다음 유한 자유 분해를 택할 수 있다.
$$
0 \to R^{\oplus n_e} \to R^{\oplus n_{e-1}} \to \ldots \to
R^{\oplus n_1} \to R \to R/I \to 0
$$
그러면 $\varphi_1 : R^{\oplus n_1} \to R$의 계수는 $1$이고,
Algebra, Proposition \ref{algebra-proposition-what-exact}에 의해
$I$가 비영인자 $y$를 포함함을 알 수 있다.
$\mathfrak p_1, \ldots, \mathfrak p_n$을 $R$의 연관 소아이디얼들이라 하자.
Algebra, Lemma \ref{algebra-lemma-finite-ass}를 보라.
$J/I^2 = F$인 아이디얼 $I^2 \subset J \subset I$를 택하자.
$y^2 \in J$이고 $y^2 \not \in \mathfrak p_i$이므로 모든 $i$에 대해
$J \not \subset \mathfrak p_i$이다. Algebra, Lemma
\ref{algebra-lemma-ass-zero-divisors}를 보라. Nakayama의 Lemma
(Algebra, Lemma \ref{algebra-lemma-NAK})에 의해
$J \not \subset \mathfrak m J + I^2$이다. Algebra, Lemma
\ref{algebra-lemma-silly}에 의해 $x \in J$를
$x \not \in \mathfrak m J + I^2$이고 $i = 1, \ldots, n$에 대해
$x \not \in \mathfrak p_i$가 되도록 택할 수 있다. 그러면 $x$는
비영인자이고, $x$의 $I/I^2$에서의 상은 Nakayama의 Lemma에 의해
직합인자 $J/I^2 \cong R/I$를 생성한다.
\end{proof}

\begin{lemma}
\label{dpa-lemma-vasconcelos}
\begin{reference}
\cite[Theorem 1.1]{Vasconcelos}의 국소 형태
\end{reference}
$R$을 뇌터 국소환이라 하자. $I \subset R$을 $R$ 위에서 유한 사영
차원을 갖는 아이디얼이라 하자. $F \subset I/I^2$가 계수 $r$인 자유
직합인자이면, $x_1 \bmod I^2, \ldots, x_r \bmod I^2$가 $F$를 생성하는
정칙열 $x_1, \ldots, x_r \in I$가 존재한다.
\end{lemma}

\begin{proof}
$r = 0$이면 증명할 것이 없다. $r > 0$이라고 가정하자. Lemma
\ref{dpa-lemma-free-summand-in-ideal-finite-proj-dim}에 의해 $x \in I$를,
$x$가 비영인자이고 $x \bmod I^2$가 $F$의 $R/I$와 동형인 한 직합인자를
생성하도록 택할 수 있다. 환 $R' = R/(x)$과 아이디얼 $I' = I/(x)$을
생각하자. 물론 $R'/I' = R/I$이다. 짧은 완전열
$$
0 \to R/I \xrightarrow{x} I/xI \to I' \to 0
$$
은 분할된다. 사상 $I/xI \to I/I^2$가 $xR/xI$를 한 직합인자로
보내기 때문이다. 이제 $I/xI = I \otimes_R^\mathbf{L} R'$은 $R'$ 위에서
유한 사영 차원을 갖는다. More on Algebra, Lemmas
\ref{more-algebra-lemma-perfect-module}와
\ref{more-algebra-lemma-pull-perfect}를 보라. 따라서 직합인자 $I'$도
$R'$ 위에서 유한 사영 차원을 갖는다. 한편 짧은 완전열
$0 \to xR/xI \to I/I^2 \to I'/(I')^2 \to 0$이 있으므로,
$I'/(I')^2$는 계수 $r - 1$인 자유 직합인자
$F' = F/(R/I \cdot x)$를 갖는다.
% P03-STACKS-E1029
$r$에 대한 귀납법으로 그 합동류들이 $F'$를 자유롭게 생성하는 정칙열
$x'_2, \ldots, x'_r \in I'$을 택할 수 있다.
% P03-STACKS-E1030
$x_1 = x$이고 $x_2, \ldots, x_r$가 $R$에서
$x'_1, \ldots, x'_r$을 올리는 임의의 원소들이면 Lemma가 성립한다.
\end{proof}

\begin{proposition}
\label{dpa-proposition-regular-ideal}
\begin{reference}
\cite[Corollary 1]{Vasconcelos}의 한 변형. 또한
\cite{Iyengar}와 \cite{Ferrand-lci}를 보라.
\end{reference}
$R$을 뇌터 환이라 하자. $I \subset R$을 유한 사영 차원을 갖고
$I/I^2$가 $R/I$ 위에서 유한 국소 자유인 아이디얼이라 하자. 그러면
$I$는 정칙 아이디얼이다(More on Algebra, Definition
\ref{more-algebra-definition-regular-ideal}).
\end{proposition}

\begin{proof}
Algebra, Lemma \ref{algebra-lemma-regular-sequence-in-neighbourhood}에 의해
모든 $\mathfrak p \supset I$에 대해 $I_\mathfrak p \subset R_\mathfrak p$가
정칙열로 생성됨을 보이면 충분하다. 따라서 $R$이 국소환이라고 가정해도
된다. $I/I^2$의 계수가 $r$이면 Lemma \ref{dpa-lemma-vasconcelos}에 의해
$I/I^2$를 생성하는 정칙열 $x_1, \ldots, x_r \in I$를 얻는다.
Nakayama(Algebra, Lemma \ref{algebra-lemma-NAK})에 의해 $I$가
$x_1, \ldots, x_r$로 생성됨을 알 수 있다.
\end{proof}

\noindent
환들의 임의의 국소 완전교차 준동형 $A \to B$에 대해 소박한 여접 복합체
$\NL_{B/A}$는 $[-1, 0]$ 안에서 Tor-진폭을 갖는 완전 복합체이다.
More on Algebra, Lemma \ref{more-algebra-lemma-lci-NL}를 보라.
위의 결과를 사용하면 이것이 어떤 경우에는 국소 완전교차 준동형을
특징지음을 보일 수 있다.

\begin{lemma}
\label{dpa-lemma-perfect-NL-lci}
$A \to B$를 완전한(More on Algebra, Definition
\ref{more-algebra-definition-pseudo-coherent-perfect})
뇌터 환들의 환 준동형이라 하자. 다음은 동치이다.
\begin{enumerate}
\item $\NL_{B/A}$는 $[-1, 0]$ 안에서 Tor-진폭을 갖는다,
\item $\NL_{B/A}$는 $[-1, 0]$ 안에서 Tor-진폭을 갖는 $D(B)$의 완전
대상이다,
\item $A \to B$는 국소 완전교차이다(More on Algebra, Definition
\ref{more-algebra-definition-local-complete-intersection}).
\end{enumerate}
\end{lemma}

\begin{proof}
$B = A[x_1, \ldots, x_n]/I$로 쓰자. 그러면 $\NL_{B/A}$는 다음
$B$-가군들의 복합체로 표현된다.
$$
I/I^2 \longrightarrow \bigoplus B \text{d}x_i
$$
여기서 $I/I^2$는 차수 $-1$에 놓인다. 차수 $0$의 항이 유한 자유이므로,
이 복합체가 $[-1, 0]$ 안에서 Tor-진폭을 가질 필요충분조건은
$I/I^2$가 평탄 $B$-가군인 것이다. More on Algebra, Lemma
\ref{more-algebra-lemma-last-one-flat}를 보라. $I/I^2$가 유한 $B$-가군이고
$B$가 뇌터 환이므로, 이는 $I/I^2$가 유한 국소 자유 $B$-가군일
필요충분조건이기도 하다(Algebra, Lemma
\ref{algebra-lemma-finite-projective}). 따라서 (1)과 (2)의 동치는 명백하다.
또한 Proposition \ref{dpa-proposition-regular-ideal}을 적용하면 (1)과 (3)의
동치도 따른다(정칙 아이디얼은 Koszul 정칙 아이디얼인 동시에
준정칙 아이디얼임을 관찰하라. More on Algebra, Section
\ref{more-algebra-section-ideals}를 보라).
\end{proof}

\begin{lemma}
\label{dpa-lemma-flat-fp-NL-lci}
$A \to B$를 유한 표시인 평탄 환 준동형이라 하자. 다음은 동치이다.
\begin{enumerate}
\item $\NL_{B/A}$는 $[-1, 0]$ 안에서 Tor-진폭을 갖는다,
\item $\NL_{B/A}$는 $[-1, 0]$ 안에서 Tor-진폭을 갖는 $D(B)$의 완전
대상이다,
\item $A \to B$는 신토믹이다(Algebra, Definition
\ref{algebra-definition-lci}),
\item $A \to B$는 국소 완전교차이다(More on Algebra, Definition
\ref{more-algebra-definition-local-complete-intersection}).
\end{enumerate}
\end{lemma}

\begin{proof}
(3)과 (4)의 동치는 More on Algebra, Lemma
\ref{more-algebra-lemma-syntomic-lci}이다.

\medskip\noindent
$A \to B$가 신토믹이면 다음 밂 도식을 찾을 수 있다.
$$
\xymatrix{
B_0 \ar[r] & B \\
A_0 \ar[r] \ar[u] & A \ar[u]
}
$$
여기서 $A_0 \to B_0$는 신토믹이고 $A_0$는 뇌터 환이다. Algebra,
Lemmas \ref{algebra-lemma-limit-module-finite-presentation}와
\ref{algebra-lemma-colimit-lci}를 보라. Lemma
\ref{dpa-lemma-perfect-NL-lci}에 의해 $\NL_{B_0/A_0}$는 $[-1, 0]$ 안에서
Tor-진폭을 갖는 완전 대상이다. More on Algebra, Lemma
\ref{more-algebra-lemma-base-change-NL-flat}에 의해
$\NL_{B/A} = \NL_{B_0/A_0} \otimes_{B_0}^\mathbf{L} B$에도 같은 사실이
성립함을 알 수 있다(More on Algebra, Lemmas
\ref{more-algebra-lemma-pull-tor-amplitude}와
\ref{more-algebra-lemma-pull-perfect}도 보라). 이는 (3)이 (2)를 함의함을
증명한다.

\medskip\noindent
(1)을 가정하자. More on Algebra, Lemma
\ref{more-algebra-lemma-base-change-NL-flat}
에 의해 $k$가 체인 모든 환 준동형 $A \to k$에 대해
$\NL_{B \otimes_A k/k}$가 $[-1, 0]$ 안에서 Tor-진폭을 가짐을 알 수 있다
(More on Algebra, Lemma \ref{more-algebra-lemma-pull-tor-amplitude}를 보라).
따라서 Lemma \ref{dpa-lemma-perfect-NL-lci}에 의해
$k \to B \otimes_A k$는 국소 완전교차 준동형이다. 그러므로 정의상
$A \to B$는 신토믹이다. 이는 (1)이 (3)을 함의함을 증명하며,
증명이 끝난다.
\end{proof}





\section{Koszul 복합체와 Tate 분해}
\label{dpa-section-koszul-vs-tate}

\noindent
이 절에서는 More on Algebra, Lemma
\ref{more-algebra-lemma-sequence-Koszul-complexes}의 결과를 미분 등급
구조와 양립하는 분할 거듭제곱이 주어진 미분 등급 대수의 범주로
``올린다''(이 절에서는 Koszul 복합체를 사슬 복합체로 나타내지만,
인용된 곳에서는 공사슬 복합체를 사용함에 유의하라).

\medskip\noindent
$R$을 환이라 하자. $I \subset R$을 $f_1, \ldots, f_r \in R$로 생성되는
아이디얼이라 하자. $n \geq 1$에 대해
$$
K_n = K_{n, \bullet} = R\langle \xi_1, \ldots, \xi_r\rangle
$$
을 $\xi_i$가 차수 $1$에 있고 $\text{d}(\xi_i) = f_i^n$인 미분 등급
Koszul 대수라 하자. 그 위에는 유일한 분할 거듭제곱 구조가 존재한다
(Definition \ref{dpa-definition-divided-powers-dga}에서와 같다).
Example \ref{dpa-example-adjoining-odd}를 보라. $m > n$에 대해 전이 사상
$K_m \to K_n$은 $\xi_i$를 $f_i^{m - n}\xi_i$로 보내는, 분할
거듭제곱과 양립하는 미분 등급 대수 사상이다.

\begin{lemma}
\label{dpa-lemma-lift-tate-to-koszul}
위의 상황에서 $R$이 뇌터 환이면, 모든 $n$에 대해 $N \geq n$과 사상들
$$
K_N \to A \to R/(f_1^N, \ldots, f_r^N)\quad\text{and}\quad A \to K_n
$$
이 존재하여 다음 성질들을 만족한다.
\begin{enumerate}
\item $(A, \text{d}, \gamma)$는 Definition
\ref{dpa-definition-divided-powers-dga}에서와 같다,
\item $A \to R/(f_1^N, \ldots, f_r^N)$은 유사동형이다,
\item 합성 $K_N \to A \to R/(f_1^N, \ldots, f_r^N)$은 표준 사상이다,
\item 합성 $K_N \to A \to K_n$은 전이 사상이다,
\item $A_0 = R \to R/(f_1^N, \ldots, f_r^N)$은 표준 전사이다,
\item $A$는 각 차수마다 유한 개의 생성원을 갖는 $R$ 위의 등급 분할
거듭제곱 다항식 대수이다,
\item $A \to K_n$은 분할 거듭제곱과 양립하는 미분 등급 $R$-대수의
준동형이며, 차수 $0$의 호몰로지에서 표준 사상
$R/(f_1^N, \ldots, f_r^N) \to R/(f_1^n, \ldots, f_r^n)$을 유도한다.
\end{enumerate}
% P03-STACKS-E598
조건 (4)는 Example \ref{dpa-example-adjoining-odd} 또는
\ref{dpa-example-adjoining-even}에서와 같이 각 차수 $> 0$마다 유한한 변수
집합 $T$를 차례로 첨가하여 $A_0$로부터 $A$를 구성한다는 뜻이다.
\end{lemma}

\begin{proof}
$n$을 고정하자. $r = 0$이면 단순히 $A = R$로 택하면 된다.
$r > 0$이라고 가정하자. More on Algebra, Lemma
\ref{more-algebra-lemma-sequence-Koszul-complexes}를 사슬 복합체의 언어로
옮겨 적용하면 다음을 택할 수 있다.
$$
n_{r} > n_{r - 1} > \ldots > n_1 > n_0 = n
$$
이때 위의 Koszul 대수들 사이의 전이 사상
$K_{n_{i + 1}} \to K_{n_i}$은 차수 $> 0$의 호몰로지에서 영사상을
유도한다. $N = n_r$로 놓고 Lemma를 증명하겠다.

\medskip\noindent
Lemma \ref{dpa-lemma-tate-resolution}의 명제와 증명에서와 똑같이 $A$를
구성한다. 따라서 다음을 얻는다.
$$
A = \colim A(m),\quad\text{and}\quad
A(0) \to A(1) \to A(2) \to \ldots \to R/(f_1^N, \ldots, f_r^N)
$$
이는 성질 (1), (2), (5), (6)을 곧바로 준다. 증명을 마치기 위해
$R$-대수 사상 $K_N \to A \to K_n$을 구성하겠다. 이를 위해 다음을
구성한다.
\begin{enumerate}
\item 동형 $A(1) \to K_N = K_{n_r}$,
\item 사상 $A(2) \to K_{n_{r - 1}}$,
\item $\ldots$
\item 사상 $A(r) \to K_{n_1}$,
\item 사상 $A(r + 1) \to K_{n_0} = K_n$,
\item 사상 $A \to K_n$.
\end{enumerate}
각 단계에서 구성되는 사상은 양립하는 분할 거듭제곱이 주어진 미분 등급
대수들 사이의 사상이다. 각 사상은 Koszul 대수들 사이의 전이 사상을
통해 앞의 사상과 양립하며, 차수 $0$의 호몰로지에서 명백한 표준 사상을
유도한다.

\medskip\noindent
$A(0) = R$임을 기억하자. $m = 1$에 대해 Lemma
\ref{dpa-lemma-tate-resolution}의 증명은 $T_i$가 차수 $1$에 있고
$\text{d}(T_i) = f_i^N$인 $A(1) = R\langle T_1, \ldots, T_r\rangle$을
택한다. 즉, $f_i^N$들은 $A(0) \to R/(f_1^N, \ldots, f_r^N)$의 핵의
생성원들이다. 따라서 $A(1) \to K_N = K_{n_r}$에는 다음 사상을 사용한다.
$$
\varphi_1 : A(1) \longrightarrow K_{n_r},\quad T_i \longmapsto \xi_i
$$
이는 동형이다.

\medskip\noindent
$m = 2$에 대해 Lemma \ref{dpa-lemma-tate-resolution}의 증명에 나오는 구성은
생성원 $e_1, \ldots, e_t \in \Ker(\text{d} : A(1)_1 \to A(1)_0)$을 택한다.
그런 다음 $T_i$가 차수 $2$에 있고 $\text{d}(T_i) = e_i$인 분할
거듭제곱 다항식 대수
$A(2) = A(1)\langle T_1, \ldots, T_t\rangle$을 취한다.
$\varphi_1(e_i)$가 $K_{n_r}$의 순환원이므로, 위에서 $n_r$과
$n_{r - 1}$을 택한 방식에 의해 그 $K_{n_{r - 1}}$에서의 상은
% P03-STACKS-E1031
쌍대경계이다. 따라서 다음 가환 도식을 구성할 수 있다.
$$
\xymatrix{
A(1) \ar[d] \ar[r]_{\varphi_1} & K_{n_r} \ar[d] \\
A(2) \ar[r]^{\varphi_2} & K_{n_{r - 1}}
}
$$
$T_i$를 그 경계가 $\varphi_1(e_i)$의 상인 차수 $2$의 원소로 보내면 된다.
% P03-STACKS-E599
분할 거듭제곱 다항식 대수의 보편성에 의해 사상 $\varphi_2$가 존재하며,
이는 미분 및 분할 거듭제곱과 양립한다.

\medskip\noindent
$m = 2, \ldots, r$에 대해 대수 $A(m)$과 사상
$\varphi_m : A(m) \to K_{n_{r + 1 - m}}$을 똑같은 방식으로 구성한다.

\medskip\noindent
사상 $A(r) \to K_{n_1}$이 주어지면 합성
$H_r(A(r)) \to H_r(K_{n_1}) \to H_r(K_{n_0}) \subset (K_{n_0})_r$이
영임을 알 수 있다. 따라서 $A(r + 1)$의 새로운 다항식 생성원들을
영으로 보내어 이를 $A(r + 1) \to K_{n_0} = K_n$으로 연장할 수 있다.

\medskip\noindent
$A(r + 1) \to K_{n_0} = K_n$을 구성한 뒤에는 새로운 다항식 생성원들을
영으로 보내는 유일하게 가능한 방식으로 $A(r + 2), A(r + 3), \ldots$에
연장하면 된다. 이로써 증명이 끝난다.
\end{proof}

\begin{remark}
\label{dpa-remark-pro-system-koszul}
위의 상황에서 $R$이 뇌터 환이면 정수열
$1 = n_0 < n_1 < n_2 < \ldots $을 귀납적으로 택하여
$i = 1, 2, 3, \ldots$에 대해 Lemma
\ref{dpa-lemma-lift-tate-to-koszul}에서와 같은 사상들
$K_{n_i} \to A_i \to R/(f_1^{n_i}, \ldots, f_r^{n_i})$
과 $A_i \to K_{n_{i - 1}}$이 존재하도록 할 수 있다.
% P03-STACKS-E1032
합성 $A_{i + 1} \to K_{n_i} \to A_i$를 $A_{i + 1} \to A_i$로 나타내자.
그러면 도식
$$
\xymatrix{
K_{n_1} \ar[d] &
K_{n_2} \ar[d] \ar[l] &
K_{n_3} \ar[d] \ar[l] &
\ldots \ar[l] \\
A_1 \ar[d] &
A_2 \ar[l] \ar[d] &
A_3 \ar[l] \ar[d] &
\ldots \ar[l] \\
K_1 &
K_{n_1} \ar[l] &
K_{n_2} \ar[l] &
\ldots \ar[l]
}
$$
은 가환한다. 이로부터 양립하는 분할 거듭제곱을 갖는 미분 등급
$R$-대수의 범주에서 역계 $(K_n)$과 $(A_n)$이 pro-동형임을 알 수 있다.
\end{remark}

\begin{lemma}
\label{dpa-lemma-compute-cohomology-adjoin-variable}
$(A, \text{d}, \gamma)$, $d \geq 1$, $f \in A_{d - 1}$,
$A\langle T \rangle$이 Lemma \ref{dpa-lemma-extend-differential}에서와 같다고 하자.
\begin{enumerate}
% P03-STACKS-E600
\item $d = 1$이면 다음 긴 완전열이 있다.
$$
\ldots \to H_0(A) \xrightarrow{f} H_0(A) \to H_0(A\langle T \rangle) \to 0
$$
\item $d = 2$이면 $H_{i + j}(A\langle T \rangle)$로 수렴하는 유계
스펙트럼 열
$(E_1)_{i, j} = H_{j - i}(A) \cdot T^{[i]}$
이 있다. 미분
$(d_1)_{i, j} : H_{j - i}(A) \cdot T^{[i]} \to
H_{j - i + 1}(A) \cdot T^{[i - 1]}$
은 $\xi \cdot T^{[i]}$를 $f \xi \cdot T^{[i - 1]}$의 류로 보낸다.
% P03-STACKS-E601
\item 필요에 따라 다른 차수의 경우를 여기에 더 추가한다.
\end{enumerate}
\end{lemma}

\begin{proof}
$d = 1$이면 다음 복합체들의 짧은 완전열이 있다.
$$
0 \to A \to A\langle T \rangle \to A \cdot T \to 0
$$
이로부터 결과 (1)이 쉽게 따른다. $d = 2$이면 $A\langle T \rangle$을
다음 부분복합체들을 갖는 여과 사슬 복합체로 본다.
$$
F^pA\langle T \rangle = \bigoplus\nolimits_{i \leq p} A \cdot T^{[i]}
$$
Homology, Section \ref{homology-section-filtered-complex}의 스펙트럼 열을
사슬 복합체로 옮겨 적용하면 (2)를 얻는다.
\end{proof}

\noindent
다음 Lemma는 뒤에서 필요하다.

\begin{lemma}
\label{dpa-lemma-construct-some-maps}
위의 상황에서 모든 $n \geq t \geq 1$에 대해 $N > n$과 사상
$$
K_t \longrightarrow K_n \otimes_R K_t
$$
이 왼쪽 미분 등급 $K_N$-가군의 유도 범주에 존재하여, 곱셈 사상과 어느
쪽 순서로 합성해도 전이 사상이 된다.
\end{lemma}

\begin{proof}
먼저 $r = 1$인 경우를 증명하자. $f = f_1$로 놓고
$K_t = R\langle x \rangle$,
$K_n = R\langle y \rangle$, $K_N = R\langle z \rangle$로 쓰자.
여기서 $x$, $y$, $z$는 차수 $1$에 있고
$\text{d}(x) = f^t$, $\text{d}(y) = f^n$,
$\text{d}(z) = f^N$이다. 모든 $N > t$에 대해 다음 유사동형이 있다고
주장한다.
$$
B_{N, t} = R\langle x, z, u \rangle
\longrightarrow
K_t,\quad
x \mapsto x,\quad
z \mapsto f^{N - t}x,\quad
u \mapsto 0
$$
여기서 왼쪽은 변수 $x$와 $z$가 차수 $1$에 있고 $u$가 차수 $2$에 있으며
$\text{d}(x) = f^t$, $\text{d}(z) = f^N$,
$\text{d}(u) = z - f^{N - t}x$인 분할 거듭제곱 다항식 대수를 뜻한다.
주장을 증명하기 위해 $H_*(R\langle x, z\rangle)$의 다음 세 부분가군이
같음을 관찰하자.
\begin{enumerate}
\item $H_*(R\langle x, z\rangle) \to H_*(K_t)$의 핵,
\item 사상
$z - f^{N - t}x : H_*(R\langle x, z\rangle) \to H_*(R\langle x, z\rangle)$
의 상,
\item 사상
$z - f^{N - t}x : H_*(R\langle x, z\rangle) \to H_*(R\langle x, z\rangle)$.
의 핵.
\end{enumerate}
이 관찰은 생략할 직접 계산으로 증명된다\footnote{힌트: 
$z' = z - f^{N - t}x$로 놓으면
$R\langle x, z\rangle = R\langle x, z'\rangle$이고 $\text{d}(z') = 0$임을
알 수 있다. 또한 사상 $R\langle x, z'\rangle \to K_t$는 $z'$을 영으로
보내는 사상이다.}. 이제 Lemma
\ref{dpa-lemma-compute-cohomology-adjoin-variable}의 (2)를 적용하면 주장이
참임을 알 수 있다.

\medskip\noindent
$z$를 $z$로 보내는 미분 등급 $R$-대수의 준동형
$K_N \to B_{N, t}$을 통해 $B_{N, t} \to K_t$를 왼쪽 미분 등급
$K_N$-가군의 유사동형으로 볼 수 있다. Lemma의 명제에 나오는 화살표를
정의하기 위해 다음 준동형을 사용한다.
$$
B_{N, t} = R\langle x, z, u \rangle \to K_n \otimes_R K_t,\quad
x \mapsto 1 \otimes x,\quad
z \mapsto f^{N - n}y \otimes 1,\quad
u \mapsto - f^{N - n - t}y \otimes x
$$
$N \geq n + t$라고 가정하면 이는 의미가 있다. 곱셈 사상과 앞이나 뒤에서
합성한 것이 전이 사상임은 간단하고 흥미로운 계산으로 확인된다.

\medskip\noindent
$r > 1$이면 각 Koszul 대수를 변수 $1$개의 Koszul 대수들의 텐서곱으로 쓰고
앞의 구성을 적용한다. 다시 말해 다음과 같이 쓴다.
$$
K_t = R\langle x_1, \ldots, x_r\rangle =
R\langle x_1\rangle \otimes_R \ldots \otimes_R R\langle x_r\rangle
$$
여기서 $x_i$는 차수 $1$에 있고 $\text{d}(x_i) = f_i^t$이다.
$r > 1$인 경우에는 다음을 사용한다.
$$
B_{N, t} =
R\langle x_1, z_1, u_1 \rangle
\otimes_R \ldots \otimes_R
R\langle x_r, z_r, u_r \rangle
$$
여기서 $x_i, z_i$는 차수 $1$에 있고 $u_i$는 차수 $2$에 있으며
$\text{d}(x_i) = f_i^t$, $\text{d}(z_i) = f_i^N$,
$\text{d}(u_i) = z_i - f_i^{N - t}x_i$이다. 텐서곱 사상
$B_{N, t} \to K_t$는 자유 $R$-가군들의 위로 유계인 복합체들 사이의
% P03-STACKS-E1033
유사동형들의 텐서곱이므로 유사동형이다. 마지막으로 사상
$$
B_{N, t}
\to
K_n \otimes_R K_t =
R\langle y_1, \ldots, y_r\rangle \otimes_R
R\langle x_1, \ldots, x_r\rangle
$$
을 $r = 1$인 경우에 구성한 사상들의 텐서곱으로 정의한다. 또는 단순히
규칙 $x_i \mapsto 1 \otimes x_i$,
$z_i \mapsto f_i^{N - n}y_i \otimes 1$,
$u_i \mapsto - f_i^{N - n - t}y_i \otimes x_i$로 정의해도 된다.
$N \geq n + t$이면 이는 의미가 있다. 세부 사항은 생략한다.
\end{proof}
